
\subsection{The ID-Algorithm}

\Joris{TODO: check for new definition of intervention node}

        Consider an IL-BCN:
            \[M= \lp G^+=\lp J,(V,U),E^+ \rp, \lp P_v ( X_v|\doit( X_{\Pa^{G^+}(v)})) \rp_{v \in V \cup U} \rp,\]
        with observational CADMG $G$.
        For subsets $A,B,C \ins V$ we want to infer the conditional interventional distribution $P(X_A|X_B,\doit(X_{J \cup C}))$ 
        in terms of (repeated products and) conditional marginals of the observational Markov kernel $P(X_V|\doit(X_J))$ 
        and knowledge of $G$. 
        In this subsection we will restrict ourselves to the case $P(X_A|\doit(X_{J \cup C}))$ with no conditioning and 
        present the \emph{ID-algorithm}, which can tell if this is possible or not (in precise terms), 
        and if so, provides us with a formula to do so. 
        We will start this subsection with a series of necessary definitions, notations and lemmata.
        The main references are \cite{Pearl09,GP95s,Tian2002,TP02s,Tian04s,SP06,HV06s,HV08s,RERS17,FM20}.

\subsubsection{Core Definitions and Notations}

\begin{Def}[Identifiability of interventional distributions/Markov kernels] 
    \label{def:id-causal-effect}
    Let $G=(J,V,E,L)$ be a CADMG and $B \ins V$ and $C \ins J \cup V$ disjoint subsets.
    We say that the interventional distribution/Markov kernel of $C$ onto $B$ is \emph{identifiable} from $G$,
    or, more in generic symbols, that $P(X_B|\doit(X_{J \cup C}))$ is \emph{identifiable} from $P(X_V|\doit(X_J))$ (and $G$),
    if for every two IL-CBNs $M_1$ and $M_2$ with the same:
    \begin{enumerate}
        \item observational CADMG $G_1=G_2=G$, and: 
        \item underlying spaces $\Xcal_{1,v}=\Xcal_{2,v}=:\Xcal_v$ for $v \in J \cup V$, and:
        \item observational Markov kernels $P_1(X_V|\doit(X_J)) =P_2(X_V|\doit(X_J))$,
    \end{enumerate}
    we also have the equality of the interventional Markov kernels:
    \[ P_1(X_B|\doit(X_{J \cup C})) = P_2(X_B|\doit(X_{J \cup C})). \]
    Sometimes we further restrict the class of CBNs to define/achieve identifiability, e.g.\ by adding
    ``for linear Gaussian CBNs'' or ``for discrete CBNs with strictly positive mass functions'', etc., and then only require $M_1$ and $M_2$ 
    to come from such classes.
\end{Def}

In the following we introduce a somewhat more vague, but constructive, notion of identifiability, which we coin \emph{\trackability} that allows us to follow certain marginalization, conditioning and multiplication steps to arrive at the wanted interventional Markov kernel.

\begin{Def}[Trackability (up to specifications)]
    Let $G=(J,V,E,L)$ be a CADMG and $B \ins V$ and $C \ins J \cup V$ disjoint subsets.
    \begin{enumerate}
        \item We say that the interventional distribution/Markov kernel of $C$ onto $B$ is \emph{\trackable} from $G$,
            or simply that $P(X_B|\doit(X_{J \cup C}))$ is \emph{\trackable} from $P(X_V|\doit(X_J))$ (and $G$),
    if there exists a finite sequence of operations, only involving marginalization, conditioning and multiplication of Markov kernels,
    applied to previously determined Markov kernels, starting from $P(X_V|\doit(X_J))$, with predetermined target sets $T_n \ins G$,
    indicating on which variable the operation is applied to, such that for every IL-CBNs $M$ with observational CADMG $G$ 
    we can compute $P(X_B|\doit(X_{J \cup C}))$ from $P(X_V|\doit(X_J))$ when we follow the above sequence of operations (and this should work no matter which version of conditional Markov kernels were used).
\item We say that $P(X_B|\doit(X_{J \cup C}))$ is \emph{\trackable\ up to specifications} ``xyz'' from $P(X_V|\doit(X_J))$ (and $G$), if the same as above holds true, but whenever we condition,  which leads to a Markov kernel only up so some null sets, we use the specifications ``xyz'' to pick a certain version of conditional Markov kernel at each step such that following the pre-specified operations leads to $P(X_B|\doit(X_{J \cup C}))$. 
\item We say that $P(X_B|\doit(X_{J \cup C}))$ is \emph{\trackable\ up to oracle choices} from $P(X_V|\doit(X_J))$ (and $G$) if there exists a conditional Markov kernel at each conditioning step (``chosen by an oracle that knows $M$'') 
    such that following these operations leads to $P(X_B|\doit(X_{J \cup C}))$. 
\item    Again, we sometimes further restrict the class of CBNs to define/achieve \trackability\ (up to oracle choices), 
    e.g.\ by adding ``for linear Gaussian CBNs'' or ``for discrete CBNs with strictly positive mass functions'', etc., 
    and then only allow $M$ to come from such classes.
    \end{enumerate}
\end{Def}


\begin{Eg}
    To illustrate how such a series of operations could look like
    consider the CADMG $G$ from Figure \ref{fig:dag-3} with $V=\lC v_1,v_2,v_3 \rC$. 
    We assume that the observational Markov kernel $P(X_1,X_2,X_3)$ is given.
    A list of operations could look like:
    \begin{enumerate}
        \item Condition $P(X_1,X_2,X_3)$ on $(X_1,X_2)$ and get a version $P(X_3|X_1,X_2)$. 
        \item[] Further specification could be (if possible): ``take a continuous version'' or ``take a version that is only dependent on variables $X_1$'' or ''take a strictly positive version''.
        \item Marginalize out $(X_2,X_3)$ from $P(X_1,X_2,X_3)$ and get $P(X_1)$.
        \item Take the product of the previous two Markov kernels: $P(X_3|X_1,X_2) \otimes P(X_1)$.
        \item Marginalize out $X_1$ from the last Markov kernel and get: $P(X_3|X_1,X_2) \circ P(X_1)$.
    \end{enumerate}
\end{Eg}

\begin{Lem}
    Let $G=(J,V,E,L)$ be a CADMG.
    \begin{enumerate} 
        \item If $P(X_B|\doit(X_{J \cup C}))$ is \trackable\ from $P(X_V|\doit(X_J))$ then it is also identifiable.
        \item If $P(X_B|\doit(X_{J \cup C}))$ is \trackable\ up to oracle choices from $P(X_V|\doit(X_J))$ 
            then it is also trackable (and thus identifiable) for discrete CBNs $M$ with strictly positive mass functions: 
    $p(x_V|\doit(x_J))>0$ for all $x_V$, $x_J$.
    \end{enumerate}
    \begin{proof}
        The first point is clear as the sequence of operations always ends in the same result.
        For discrete CBNs $M$ with strictly positive mass functions conditional Markov kernels are unambiguous, thus a sequence
        of marginalization, conditioning and products always leads to the same result.
        Note that marginals, conditionals and products of strictly positive mass functions also are strictly positive mass functions.
    \end{proof}
\end{Lem}



\begin{figure}[ht]
\centering
\begin{tikzpicture}[scale=1, transform shape]
  \begin{scope}
    \node[ndout] (v1) at (0,0) {$v_1$};
    \node[ndout] (v2) at (2,0) {$v_2$};
    \draw[arout] (v1) to (v2);
    \node (a) at (-1,0.5) {(a)};
  \end{scope}
  \begin{scope}[xshift=5cm]
    \node[ndout] (v1) at (0,0) {$v_1$};
    \node[ndout] (v2) at (2,0) {$v_2$};
    \draw[arout] (v1) to (v2);
    \draw[arlat, bend left] (v1) to (v2);
    \node (a) at (-1,0.5) {(b)};
  \end{scope}
  \begin{scope}[xshift=10cm]
    \node[ndout] (v1) at (0,0) {$v_1$};
    \node[ndout] (v3) at (2,0) {$v_{3}$};
    \node[ndout] (v2) at (4,0) {$v_2$};
    \draw[arout] (v1) to (v3);
    \draw[arout] (v3) to (v2);
    \draw[arlat, bend left] (v1) to (v2);
    \node (a) at (-1,0.5) {(c)};
  \end{scope}
\end{tikzpicture}
\caption{(a) A DAG with two nodes. (b) An ADMG with two nodes. (c) An ADMG with three nodes.
The interventional distribution $P(X_{2}|\doit(X_{1}))$ is trackable up to oracle choices from $P(X_{1},X_{2},X_{3})$ in (a) and (c), but not in (b).}
\label{fig:dag-2}
\end{figure}

\begin{Eg}
    Consider the DAG $G=(V,E)$ from Figure \ref{fig:dag-2} (a) with $V=\lC v_1,v_2 \rC$ and $E=\lC v_1\tuh v_2 \rC$.
    Let $\Xcal_{1}:=\lC a,b,c\rC$ and $\Xcal_{2}:=\lC 0,1 \rC$.
    Define the following Markov kernels $P_1(X_{1}) := P_1(X_{2}):=P(X_{1})$ via:
    \begin{align*}
        P(X_{1}=a) &:= \frac{1}{2}, &P(X_{1}=b) &:= \frac{1}{2}, & P(X_{1}=c) &:= 0,
    \end{align*}
    and further:
   \begin{align*}
        P_1(X_{2}|\doit(X_{1}=a)) & =  \mathrm{Bern}(1/4), & P_2(X_{2}|\doit(X_{1}=a)) & =  \mathrm{Bern}(1/4), \\
        P_1(X_{2}|\doit(X_{1}=b)) & =  \mathrm{Bern}(3/4), & P_2(X_{2}|\doit(X_{1}=b)) & =  \mathrm{Bern}(3/4), \\
        P_1(X_{2}|\doit(X_{1}=c)) & =  \mathrm{Bern}(1/8), & P_2(X_{2}|\doit(X_{1}=c)) & =  \mathrm{Bern}(7/8).
    \end{align*}
    Then we have two CBNs with observational DAG $G$, the same underlying spaces 
    and the same observational Markov kernel: 
    \[ P(X_{1},X_{2}) := P_1(X_{2}|\doit(X_{1})) \otimes P_1(X_{1}) = P_2(X_{2}|\doit(X_{1})) \otimes P_2(X_{1}), \]
    given by:
\begin{align*} M_1:=\lp G, \lp P_1(X_{1}),P_1(X_{2}|\doit(X_{1})) \rp \rp, \\
               M_2:=\lp G, \lp P_2(X_{1}),P_2(X_{2}|\doit(X_{1})) \rp \rp. 
\end{align*}
    Furthermore, consider the Markov kernel $P(X_{2}|X_{1})$ given by:
   \begin{align*}
        P(X_{2}|X_{1}=a) & =  \mathrm{Bern}(1/4),\\
        P(X_{2}|X_{1}=b) & =  \mathrm{Bern}(3/4), \\
        P(X_{2}|X_{1}=c) & =  \mathrm{Bern}(5/8).
    \end{align*}
    We thus have three different versions of the conditional Markov kernels of $P(X_{1},X_{2})$: 
 \[ P(X_{2}|X_{1}) \neq P_1(X_{2}|\doit(X_{1})) \neq P_2(X_{2}|\doit(X_{1})) \neq P(X_{2}|X_{1}). \]
 This shows that the interventional distribution $P(X_{2}|\doit(X_{1}))$ is not identifiable (and thus not \trackable) from $P(X_{1},X_{2})$ and  $G$.
 However, it is  \trackable\ up to oracle choices from $P(X_{1},X_{2})$ and $G$.
 It would thus be trackable (and identifiable) for discrete CBNs with strictly positive mass functions, e.g.\ here if
 we also put positive mass on $P(X_{1}=c) > 0$.
\end{Eg}


\begin{Not}
    Let $G=(J,V,E,L)$  be a CADMG, $<$ a topological order of $G$ and let $v \in C \ins V$.
    We then put:
    \begin{align*}
        \Anc^{[C]}(v) &:= \Anc^{G_{\doit(C^\cmpl)}}(v) \cap C, \\
        \Pred^{[C]}_<(v)&:= \Pred_<^{G_{\doit(C^\cmpl)}}(v) \cap C &&= \Pred_<^{G}(v) \cap C, \\
        \Dist^{[C]}(v)&:= \Dist^{G_{\doit(C^\cmpl)}}(v) \cap C &&= \Dist^{G_{\doit(C^\cmpl)}}(v) , \\
        \Dist^{[C]}_<(v)&:= \Dist^{[C]}(v) \cap \Pred^{[C]}_<(v), \\
        \Dst[C] &:= \lC \Dist^{[C]}(v) \st v \in C \rC.
    \end{align*}
    Similarly, if we use the subscript $\le$ we then also include $v$.
    Also note that the dependence on $G$ in the above constructions is implicit.
\end{Not}



\begin{Not}[The key interventional Markov kernels]
     Let $M$ be an IL-CBN with with observational CADMG $G=(J,V,E,L)$ and $C \ins V$ any subset.
     We will abbreviate:
    \begin{align*} \Qf[C] 
        &:= P(X_C|\doit(X_{J \cup V \sm C})) 
        = P(X_C|\doit(\cancel{X_{(J \cup V) \sm (\Pa^G(C) \cup C)}}, X_{\Pa^G(C)\sm C})),
    \end{align*}
    where the latter identification comes from Lemma \ref{lem:C-intPaC} (using the global Markov property), which is 
    ``surely'' determined, and not just up to some null-set.
    
    Note that we have the corner cases:
    \[ \Qf[V]=P(X_V|\doit(X_J)), \qquad \Qf[\emptyset]=\delta_\ast.  \]
\end{Not}


%\paragraph{Proofs - The Key Interventional Markov Kernels}\hfill

\begin{Lem}
    \label{lem:C-intPaC}
    Let $M$ be an IL-CBN with with observational CADMG $G=(J,V,E,L)$ and $C \ins V$ any subset.
    Then we have the identification:
    \[ P(X_C|\doit(X_{J \cup V \sm C})) = P(X_C|\doit(\cancel{X_{(J \cup V) \sm (\Pa^G(C) \cup C)}}, X_{\Pa^G(C)\sm C})). \]
    \begin{proof}
        This follows from the global Markov property with:
        \[ C \Perp^d_{G_{\doit(I_{V \sm (C \cup \Pa^G(C))},\Pa^G(C)\sm C)}} I_{V \sm (C \cup \Pa^G(C))} \given \Pa^G(C)\sm C.\]
        To elaborate the latter, let $P:=\Pa^G(C)\sm C$ and $W := V \sm (C \cup \Pa^G(C))$ and:
        \[ V':=V \sm P= C \dcup W, \qquad J':= (J\sm P) \dcup P, \qquad G':=G_{\doit(I_W,P)}.\]
        Now consider a walk from a node $c \in C$ to a node $j \in J' \dcup I_W$ in $G'$:
        \[\pi:\qquad c \sus \cdots \sus j.\]
        If $j \in P$ then the walk is blocked by $P$ at the endnode $j \in P$.
        So lets assume the case $j \notin P$. Then the walk is of the form:
        \[\pi:\qquad c \sus \cdots \sus w \hut j,\]
        with a $w \in W$. So we can write it further as:
        \[\pi:\qquad c =c_0 \sus c_1 \sus \cdots c_k \sus v \sus \cdots \sus w \hut j,\]
        with $c_0,\dots,c_k \in C$ for some $k \ge 0$, and $v \notin C$, the first occuring node not in $C$ (on $\pi$ from the left). 
        Note that $v=w$ is possible.
        If the edge $c_k \sus v$ is of the form $c_k \hut v$ then $v \in P$ and the walk is blocked by $P$ at the non-collider $v$.
        So we can assume the case where the edge is of the form $c_k \suh v$. 
        This means that on the subwalk $c_k \suh v \sus \cdots \sus w \hut j$ we must have at least one collider.
        This collider is then blocked by $P$ as no collider can be an ancestor of a node in $P$ inside $G'$, because $P$ consists only of input nodes of $G'$. 

        This shows the claim:
        \[ C \Perp^d_{G'} I_W \given P.\]
        %Then $v \notin \lp C \cup J \cup P \cup I_W \rp$, thus $v \in W$. 
        The rest then follows from the global Markov property.
    \end{proof}
\end{Lem}


\subsubsection{The Interventional Ordered Local Markov Property}

One of the ingredient for the ID-algorithm is the ability to \track\ (up to oracle choices) 
the interventional Markov kernel $\Qf[D]$  for districts $D$ of $G$ from $\Qf[V]$.
The key ingredient to achieve this is the \emph{interventional ordered local Markov property}, which
provides us with certain well-behaved Markov kernels that appear in factorizations of both $\Qf[V]$ and $\Qf[D]$.

\begin{Def}[The preceding Markov blanket of a node]
    Let $G=(J,V,E,L)$ be a CADMG and $<$ a topological order.
    For $v \in  V$ we make the following abbreviations:
    \begin{align*}
        G_\le(v) &:= \Pred^G_\le(v), \\% &  [A]_\le(v) &:= \Pred^{[A]}_\le(v), \\
        \Di^G_\le(v) &:= \Dist^{G_\le(v)}(v),\\% &\Di^{[A]}_\le(v) &:= \Dist^{[A]_\le(v)}(v), \\
        \Di^G_<(v) &:= \Di^G_\le(v) \sm \lC v \rC, \\%& \Di^{[A]}_<(v) &:= \Di^{[A]}_\le(v) \sm \lC v \rC, \\
        \PaD^G_<(v) &:= \Pa^G(\Di^G_\le(v)) \sm \Di^G_\le(v), \\%& \PaD^{[A]}_<(v) &:= \Pa^G(\Di^{[A]}_\le(v)) \sm \Di^{[A]}_\le(v),\\
        \MBl^G_<(v) &:= \PaD^G_<(v) \dcup \Di^G_<(v) \\ %, & \MBl^{[A]}_<(v) &:= \PaD^{[A]}_<(v) \dcup \Di^{[A]}_<(v).
                    &= \Pa^G(\Dist^{G_\le(v)}(v)) \cup \Dist^{G_\le(v)}(v) \sm \lC v \rC.
    \end{align*}
    We call $\Di^G_\le(v)$ the \emph{preceding district} and
    $\MBl^G_<(v)$ the  \emph{preceding Markov blanket} of $v$ in $G$ w.r.t.\ $<$.
    Note that this definition depends on the topological order $<$ and that we have the inclusions: 
    \[\Di^G_<(v) \ins  \MBl^G_<(v) \ins  \Pred_<^G(v).\] 
\end{Def}



\begin{Prp}[Interventional ordered local Markov property]
    \label{prp:ioLMP}
    Let $M$ be an IL-CBN with with observational CADMG $G=(J,V,E,L)$ and a fixed topological order $<$.
    Then for every $v \in V$ we have the conditional independence:
    \[ X_v \Indep_{P(X_V|\doit(X_{I_{V \sm \Di^G_\le(v)}},X_J))}  X_{\Pred^G_<(v)} \given X_{\MBl^G_<(v)}.\]
    %\[ X_v \Indep_{P(X_V|\doit(X_J))}  X_{\Pred^G_<(v)} \given X_{\MBl^G_<(v)}.\]
    In particular, 
    %when applied to $G_{\doit(I_{V \sm \Dist^G(v)})}$, 
    there exists a Markov kernel, denoted by: 
    \[Q(X_v|X_{\MBl^G_<(v)}) \qquad \text{or} \qquad Q(X_v|X_{\Di^G_<(v)},\doit(X_{\PaD^G_<(v)})), \] 
    that simultaneously is a version of:
    \[ P(X_v|X_{\Pred^{[V]}_<(v)},\doit(X_J)) \qquad \text{and} \qquad P(X_v|X_{\Pred^{[D]}_<(v)},\doit(X_{J \cup V\sm D})),\]
    for every subset $D \ins V$ with $\Di^G_\le(v) \ins D$, e.g.\ $D=\Dist^G(v)$.
    \begin{proof}
        This follows from the global Markov property, Theorem \ref{thm-gmp-mI-CBN}, together with the d-separation statement:
        \[ \lC v\rC \Perp^d_{ G_{\doit( I_{V \sm \Di^G_\le(v)} )}}  \Pred^G_<(v) \given \MBl^G_<(v).\]
        See Lemma \ref{lem:id-d-sep-3} and Lemma \ref{lem:int-dist-per-node}.
    \end{proof}
\end{Prp}

\paragraph{Proofs - The Interventional Ordered Local Markov Property}\hfill

For the next two Lemmata we introduce some shorter notations:

\begin{Not}
    Let $G=(J,V,E,L)$ be a CADMG, $<$ a topological order for $G$ and $v \in V$ a fixed node.
    For a subset $W \ins J \cup V$ we abbreviate:
    \[ W_<:=\lC w \in W \st w < v \rC, \qquad W_\le:=\lC w \in W \st w < v \,\lor\, w=v \rC, \]
    and $W_>$ and $W_\ge$ accordingly.
\end{Not}

The next Lemma is the graphical center piece that makes the ID-algorithm possible.
It could be called the graphical version of an ``interventional ordered local Markov property'',
whose distributional counterpart is stated in the Lemma after.


\begin{Lem}
    \label{lem:id-d-sep-3}
   Let $G=(J,V,E,L)$ be a CADMG, $<$ a topological order for $G$ and $v \in V$ a fixed node.
   Let $D \ins V$ be a subset such that $v \in D$ and $\Dist^{G_\le}(D_\le)\ins D$,
   where $G_\le:=\Pred^G_\le(v)$ is the ancestral subgraph of predecessors of $v$ in $G$, e.g.\ $D=\Dist^G(v)$ or $D=\Dist^{G_\le}(v)$.
   Then we have the d-separation statement:
   \[ \lC v \rC \Perp^d_{G_{\doit( I_{V\sm D_\le} )}}  V_< \given \Pa^G(D_\le) \cup D_<.  \]
   %\[ \lC v \rC \Perp^d_{G_{\doit( I_{V\sm D_\le} )}} I_{V\sm D_\le} \dcup J \dcup V_< \given \Pa^G(D_\le) \cup D_<.  \]
   \begin{proof}
       We abbreviate:
       \[\tilde G := G_{\doit( I_{V\sm D_\le} )}, \qquad F:= \Pa^G(D_\le) \cup D_<.\]
       Assume the contrary to the claim and let $\pi$ be a shortest path from $v$ to a node $w \in J \dcup I_{V \sm D} \dcup V_<$ in the graph $\tilde G$. It is clear that $w \neq v$.

       If $w \in F$ then $\pi$ is blocked by $F$ at the endnode $w$.
       So we can assume that $w \notin F$, in particular, $w \notin D_\le$.
       So there exist $v_0,\dots,v_k \in D_\le$ for some $k\ge 0$ and $\tilde w \notin D_\le$ such that $\pi$ is of the form:
       \[\pi:\qquad v=v_0 \sus \cdots \sus v_k \sus \tilde w \sus \cdots  \sus w, \]
       where $\tilde w$ is the first node from the left that is not in $D_\le$, which exists since $w \notin D_\le$. 
       Since $v_k \in D_\le$ it is clear that $\tilde w \notin I_{V\sm D_\le}$. %even $\tilde w \notin I_{V\sm D_\le}$.  
       So $\tilde w \in J \cup V \sm D_\le$. 
      
       If the edge $v_k \sus \tilde w$ is of the form $v_k \hut \tilde w$ then $\tilde w \in \Pa^G(D_\le)$. 
       So in this case $\pi$ is blocked at the non-collider $\tilde w$ by $F$.

       So we can assume the case $v_k \suh \tilde w$. This implies that $\tilde w$ cannot lie in the set of input nodes of $\tilde G$, 
       which implies that $\tilde w \notin J \cup I_{V\sm D_\le}$. With this we then get that $\tilde w \in V \sm D_\le$. 

       Assume the case that $\tilde w \in V_>$ and $\pi$ is of the form:
        \[\pi:\qquad v=v_0 \sus \cdots \sus v_k \suh \tilde w = \tilde w_0 \sus  \cdots  \sus \tilde w_m =w, \]
       where the subwalk $\tilde w_0 \sus  \cdots  \sus \tilde w_m$ has no colliders. 
       Because of the edge $v_k \suh \tilde w$ this subwalk necessarily is directed to the right and we get:
        \[\pi:\qquad v=v_0 \sus \cdots \sus v_k \suh \tilde w = \tilde w_0 \tuh  \cdots  \tuh \tilde w_m =w. \]
       Then the endnode $w$ is not an input node, $w \notin J \cup I_{V\sm D_\le}$, and $v < \tilde w \le \tilde w_m= w$, thus $w \in V_>$. 
       But this is a contradiction to $w \in J \dcup I_{V \sm D} \dcup V_<$. So this case cannot occur.

       Now assume the case that $\tilde w \in V_>$ and $\pi$ is of the form:
        \[\pi:\qquad v=v_0 \sus \cdots \sus v_k \suh \tilde w = \tilde w_0 \tuh  \cdots  \tuh \tilde w_m \hus \cdots  \sus w, \]
        for some $m \ge 0$, with a directed subwalk $\tilde w_0 \tuh  \cdots  \tuh \tilde w_m$, where $\tilde w_m$ is the first node after $\tilde w$ where a collider occurs, which could be $\tilde w$ itself. Again we have $v < \tilde w \le \tilde w_m$ and thus 
        $\tilde w_m \in V_>$. This implies that $\tilde w_m \notin \Anc^G(F)$, since $\Anc^G(F) \ins J \cup V_<$. 
        So $\pi$ is blocked by $F$ at the collider $\tilde w_m$.

       So we are left with the cases $v_k \suh \tilde w$ and $\tilde w \in V_<\sm D_\le$.
 
       Now consider the case of a directed edge $v_k \tuh \tilde w$ and $\tilde w \in V_<\sm D_\le$.
       If $v_k \neq v$ then $\pi$ is blocked at the non-collider $v_k \in D_<$ by $F$. 
       So we can assume that $v_k=v$. This implies $v < \tilde w$ and thus $\tilde w \in V_>$, which contradicts $\tilde w \in V_<$.
       %(or leads to the case already considered above). 
       So this cannot occur.

       Now consider the case of a bidirected edge $v_k \huh \tilde w$ and $\tilde w \in V_< \sm D_\le$.
       Then both nodes $v_k,\tilde w \in G_\le$ and $\tilde w \in \Dist^{G_\le}(v_k)$. By the assumption of this Lemma we have:
       \[ \tilde w \in \Dist^{G_\le}(v_k) \ins \Dist^{G_\le}(D_\le) \ins D \cap G_\le = D_\le. \]
       So $\tilde w \in D_\le$, which contradicts $\tilde w \notin D_\le$. So this case cannot neither occur.

       So we have shown that in all cases that can occur the path $\pi$ in $\tilde G$ is blocked by $F$. This shows the claim.
   \end{proof}
\end{Lem}




The last Lemma allows us to use the global Markov property for the existence of special Markov kernels that are of importance for the ID-algorithm:


\begin{Lem}[Interventional ordered local Markov property]
    \label{lem:int-dist-per-node}
    Let $M$ be an IL-CBN with with observational CADMG $G=(J,V,E,L)$ and a fixed topological order $<$ and fixed $v \in V$.
    Let $G_\le:=\Pred^G_\le(v)$ be the ancestral sub-CADMG of predecessors of $v$ in $G$ and let:
    \[ D_\le := \Dist^{G_\le}(v), \qquad F:= \Pa^G(D_\le) \cup D_\le \sm \lC v \rC. \]
   Then we have the conditional independence:
   \[ X_v \Indep_{P(X_{V_\le}|\doit(X_{I_{V\sm D_\le}},X_J))}  X_{V_<} \given X_F.\]
   In particular, there exists a Markov kernel:
   \[ Q(X_v|X_F) \] 
   that simultaneously is a version of:
   \begin{align*} 
       P(X_v|X_H,\doit(X_{J \cup V\sm S})), 
    \end{align*}
   for every subsets $H,S \ins V$ such that $D_\le \ins S \ins V$ and $F \cap S_< \ins  H \ins V_<$.
  Note that this includes the corner cases: 
  \begin{enumerate}
      \item $P(X_v|X_{V_<},\doit(X_J))$,
      \item $P(X_v|X_{F\cap V_<},\doit(X_J))$,
      \item $P(X_v|X_{S_<},\doit(X_{J \cup V\sm S}))$,
      \item $P(X_v|X_{F \cap S_<},\doit(X_{J \cup V\sm S}))$,
      \item $P(X_v|X_{D_<},\doit(X_{J \cup V\sm D_\le}))$.
  \end{enumerate}
  % \[ P(X_v|X_{V_<},\doit(X_J)), \qquad P(X_v|X_F,\doit(X_J)) \qquad \text{and} \qquad P(X_v|X_{S_<},\doit(X_{J \cup V\sm S})), \]
  % for every subset $S \ins V$ with $D_\le \ins S$.
   %Furthermore, such a $Q(X_v|X_F)$ is unique up to a $P(X_{V_<}|\doit(X_{I_{V \sm D}},X_J))$-null set.
 \begin{proof}
        By Lemma \ref{lem:id-d-sep-3} we have the d-separation:
   \[ \lC v \rC \Perp^d_{G_{\doit( I_{V\sm D_\le} )}} V_< \given F.  \]
   By the global Markov property, Theorem \ref{thm-gmp-mI-CBN}, we get: 
   \[ X_v \Indep_{P(X_{V_\le}|\doit(X_{I_{V\sm D_\le}},X_J))}  X_{V_<} \given X_F.\]
   So there exists a Markov kernel $Q(X_v|X_F)$ such that:
   \begin{align*}  P(X_{V_\le}|\doit(X_{I_{V\sm D_\le}},X_J)) = Q(X_v|X_F) \otimes P(X_{V_<}|\doit(X_{I_{V\sm D_\le}},X_J)). 
       \label{eq:q-mk-1}\tag{\#}
   \end{align*}
   For a subset $S \ins V$ with $D_\le \ins S$ we have: 
       \[ V\sm D_\le = (V\sm S) \dcup (S\sm D_\le). \] 
   %Considering $X_{I_{S \sm D_\le}} = \star$ and the two corner cases $X_{I_{V \sm S}}=\star$ and $X_{I_{V \sm S}} = x_{V \sm S}$, we get:
   By putting $X_{I_{S \sm D_\le}} = \star$ and $X_{I_{V \sm S}} = x_{V \sm S}$ we get:
   \begin{align*}
     %  P(X_{V_\le}|\doit(X_J)) &= Q(X_v|X_F) \otimes P(X_{V_<}|\doit(X_J)), \\
       P(X_{V_\le}|\doit(X_{V \sm S},X_J)) &= Q(X_v|X_F) \otimes P(X_{V_<}|\doit(X_{V \sm S},X_J)).
   \end{align*}
 Now consider another subset $H \ins V_<$ with $F \cap S_< \ins  H$. 
 Then marginalizing out $X_{V_\le \sm \lp \lC v \rC \cup H \rp}$ gives us:
   \begin{align*}
   P(X_v,X_H|\doit(X_{J \cup V \sm S})) &= Q(X_v|X_F) \otimes P(X_H|\doit(X_{J \cup V \sm S})).
   \end{align*}
 This shows that $Q(X_v|X_F)$, simultaneously, is a version of:
   \begin{align*} P(X_v|X_H,\doit(X_{J \cup V\sm S})), 
    \end{align*}
   for every subsets $H,S \ins V$ such that $D_\le \ins S \ins V$ and $F \cap S_< \ins  H \ins V_<$.
  This shows the claim.
     \end{proof}
\end{Lem}


\subsubsection{Ancestral Sets and Districts}

\begin{Lem}[Ancestral subsets are \trackable]
    \label{lem:int-anc}
    Let $M$ be an IL-CBN with with observational CADMG $G=(J,V,E,L)$ and $A \ins V$ be a subset such that 
    $A=\Anc^{[V]}(A)$.
    Then we have the equality between the interventional distribution $\Qf[A]$ and the $A$-marginal of $\Qf[V]$:
            \[ \Qf[A] = P(X_A|\doit(X_{J \cup V \sm A})) = P(X_A|\doit(X_J)). \]
    \begin{proof}
        By Lemma \ref{lem:C-intPaC} we only have to show the right identity:
        \[ \Qf[A] = P(X_A|\doit(X_{J \cup V \sm A})) \stackrel{!}{=} P(X_A|\doit(X_J)). \]
       The latter follows again from the global Markov property and the d-separation:
       \[ A \Perp^d_{G_{\doit(I_{V \sm A})}} I_{V \sm A} \given J.  \]
      % The arguments are similar to \ref{lem:C-intPaC} with the  $\Pa^G(A) \sm A \ins J$.
       This d-separation holds true since every walk from a node $v \in A$ to a node $j \in J \dcup I_{V \sm A}$ is either blocked by $J$ as the endnode $j \in J$ or is of the form:
       \[\pi:\qquad v \sus a_1 \sus \cdots a_k \sus w' \sus \cdots \sus w \hut j,\]
       with a $ w \in V \sm A$, $j \in I_{V\sm A}$, $a_1,\dots, a_k \in A$ for some $k \ge 0$ and $w' \notin A$, the first node not in $A$ on $\pi$ from the left ($w'=w$ possible).
       In case the edge $a_k \sus w'$ is of the form $a_k \hut w'$ we have:
       \[w' \in \Anc^G(A) \sm A = \Anc^G(A) \sm \Anc^{[V]}(A)  \ins J.\] 
      So in this case the walk is blocked at the non-collider $w'$ by $J$. 
       So we can consider the case where the edge is of the form $a_k \suh w'$.
       Then the subwalk $a_k \suh w' \sus \cdots \sus w \hut j$ must contain a collider. This collider can not be an ancestor of $J$, as $J$ are the input nodes. So the walk is blocked by $J$ in all cases.
    \end{proof}
\end{Lem}

\begin{Rem}[Districts are trackable up to oracle choices] 
    \label{rem:Q-dist-ident}
    Let $M$ be an IL-CBN with with observational CADMG $G=(J,V,E,L)$.
    \begin{enumerate}
        \item Since the Markov kernel $Q(X_v|X_{\MBl^G_<(v)})$ coming from the interventional ordered local Markov property, 
            see Proposition \ref{prp:ioLMP}, is a version of both:
          \[ P(X_v|X_{\Pred^{[V]}_<(v)},\doit(X_J)) \qquad \text{and} \qquad P(X_v|X_{\Pred^{[D]}_<(v)},\doit(X_{J \cup V\sm D})),\]
                for $D=\Dist^G(v)$, which are marginal conditionals of the interventional distributions $\Qf[V]$ and $\Qf[D]$, resp.,
                the $Q(X_v|X_{\MBl^G_<(v)})$'s are \trackable\ from either quantity up to oracle choices.
            \item We get the following factorization by the chain rule for every $D \in \Dst[V]$: %or unions thereof if Fa is adjusted
               \begin{align*}
                   \Qf[V] &= \bigotimes_{v \in V}^> P(X_v|X_{\Pred^{[V]}_<(v)},\doit(X_J)) &= \bigotimes_{v \in V}^> Q(X_v|X_{\MBl^G_<(v)}), \\
                   \Qf[D] &= \bigotimes_{v \in D}^> P(X_v|X_{\Pred^{[D]}_<(v)},\doit(X_{J \cup V\sm D})) &= \bigotimes_{v \in D}^> Q(X_v|X_{\MBl^G_<(v)}).
               \end{align*}
           \item In particular, $\Qf[D]$ is \trackable\ from $\Qf[V]$ up to oracle choices 
               by first determining $Q(X_v|X_{\MBl^G_<(v)})$ for $v \in D$  
               via marginalization and conditioning and then taking the product of Markov kernels in reverse order of $<$.
            \item The above seems to give us something like a factorization:
       \begin{align*}
           \Qf[V] &= \bigotimes_{v \in V}^> Q(X_v|X_{\MBl^G_<(v)}) 
            && = \lB \bigotimes_{D \in \Dst[V]}^> \bigotimes_{v \in D}^>  \rB Q(X_v|X_{\MBl^G_<(v)}) \\
            & = \lB \bigotimes_{D \in \Dst[V]}^> \rB \lp \bigotimes_{v \in D}^>  Q(X_v|X_{\MBl^G_<(v)}) \rp  
            && = \lB \bigotimes_{D \in \Dst[V]}^>\rB \Qf[D],             
       \end{align*}
         where the products in brackets are not well-defined in the naive way, 
         as the districts of a CADMG don't need to be topologically ordered.
         Nonetheless, if we are given a fixed topological order $<$ and the interventional Markov kernels $\Qf[D]$ 
         for every $D \in \Dst[V]$ 
         then we can first \track\ $Q(X_v|X_{\MBl^G_<(v)})$ from $\Qf[D]$ up to oracle choices
         for every $v \in D$ and every $D \in \Dst[V]$ and then take the product
         on the top left. 
    \end{enumerate}
\end{Rem}

\begin{Def}
    \label{not:Q-dist}
        Let $M$ be an IL-CBN with with observational CADMG $G=(J,V,E,L)$ and a fixed topological order $<$.
        Let  $\Dst \ins \Dst[V]$ be a set of districts of $G$. Then we put:
         \[ \lB \bigotimes_{D \in \Dst}^>\rB \Qf[D] := \bigotimes_{v \in \bigcup_{D \in \Dst} D}^> Q(X_v|X_{\MBl^G_<(v)}),  \]
         where the product on the right is taken in reverse topological order.
         Note that for $G':=G_{\doit(D^\cmpl)}$ and $v \in D$ we have $\MBl_<^{G'}(v)=\MBl_<^G(v)$. So the set $\MBl_<^G(v)$ 
         can be determined by the subgraph $G'$ and $<$ alone.
\end{Def}




\subsubsection{The ID-Algorithm}

Now we come to the main part of this section, the \emph{ID-algorithm} for the identification of causal effects, or, more precisely, the \trackability\ up to oracle choices of interventional Markov kernels, from the observational Markov kernel. 
The main references are \cite{Pearl09,GP95s,Tian2002,TP02s,Tian04s,SP06,HV06s,HV08s,RERS17,FM20}.

\begin{Alg}[ID-algorithm] \label{alg:id}
    Let $M$ be an IL-CBN with with observational CADMG $G=(J,V,E,L)$ and a fixed topological order $<$.
    Let $\emptyset \neq B \ins V$ and $C \ins J \cup V$ be two disjoint subsets of nodes. 
    We want to query if the interventional Markov kernel $P(X_B|\doit(X_{J \cup C}))$ is \trackable\ up to oracle choices
    from the observational Markov kernel $P(X_V|\doit(X_J))=\Qf[V]$ and $G$.
    \begin{enumerate}
        \item Put $B^C:=\Anc^{G_{\doit(C)}}(B) \sm (J \cup C) \ins V$. 
        \item[] Then $P(X_B|\doit(X_{J \cup C}))$ is the 
            $B$-marginal of $P(X_{B^C}|\doit(X_{J \cup C})) = \Qf[B^C]$.
        \item[] So we are left to determine if we can \track\ $\Qf[B^C]$ from $\Qf[V]$ up to oracle choices.
        \item Find the districts $\Dst[B^C]=\lC S_1,\dots,S_K\rC$ and put $A_{k,0}:=V$ for $k=1,\dots,K$.
        \item[] Note that $\Qf[A_{k,0}]=\Qf[V]$ is trivially \track ed from $\Qf[V]$.
        \item For each $k=1,\dots, K$ repeat the following steps recursively for $\ell \in \N$:
            \begin{enumerate}
                \item Take the district in $A_{k,\ell}$:
                        \[ D_{k,\ell} := \Dist^{[A_{k,\ell}]}(S_k) \]
                        We can \track\ $\Qf[D_{k,\ell}]$ from $\Qf[A_{k,\ell}]$ up to oracle choices by Remark \ref{rem:Q-dist-ident}.
                \item Take the ancestral closure in $D_{k,\ell}$: 
                    \[ A_{k,\ell+1}:=\Anc^{[D_{k,\ell}]}(S_k). \] 
                    We can \track\ $\Qf[A_{k,\ell+1}]$ from $\Qf[D_{k,\ell}]$ via marginalization by Lemma \ref{lem:int-anc}.

                \item If $D_{k,\ell}=A_{k,\ell}$ or $A_{k,\ell+1}=D_{k,\ell}$ then stop for this $k$ and put: 
                    \[ \idc S_k:=D_{k,\ell}.\]
                \item[] Otherwise, repeat with: $\ell \leftarrow \ell +1$.
            \end{enumerate}
                \item When the algorithm has stopped then for every $k=1,\dots,K$ we have:
                    \[ \idc S_k=\Anc^{[\idc S_k]}(S_k)=\Dist^{[\idc S_k]}(S_k) \sni S_k.\]
                    Furthermore, we have \track ed all $\Qf[\idc S_k]$'s recursively from $\Qf[V]$ up to oracle choices.
                \item If there is any $k=1,\dots,K$ with $\idc S_k \neq S_k$ then the ID-algorithm outputs: FAIL.
                \item Otherwise, we have for all $k=1,\dots,K$ that $\Qf[S_k]=\Qf[\idc S_k]$ and we can \track\
                    $\Qf[B^C]$ from $\Qf[V]$ up to oracle choices via:
                    \[ \Qf[B^C] = \lB \bigotimes_{S_k \in \Dst[B^C]}^> \rB \Qf[S_k],\]
                    and $P(X_B|\doit(X_{J \cup C}))$ as the $B$-marginal thereof.
    \end{enumerate}
\end{Alg}



\begin{Cor}[Soundness up to oracle choices]  
    \label{cor:id-alg-sound-oracle}
    The ID-algorithm \ref{alg:id} is sound up to oracle choices.
    This means that if it does not produce FAIL for input $B, C \ins G$ 
    then $P(X_B|\doit(X_{J \cup C}))$ is \trackable\ up to oracle choices from $P(X_V|\doit(X_J))$ and $G$.
    \begin{proof}
    This is clear as each step in the ID-algorithm is trackable up to oracle choices.
    Note that the operations at each step can be formulated by knowing $G$, $B$ and $C$ alone 
    without knowing $M$ in advance.
 \begin{comment}
   First note that all interventional Markov kernels $\Qf[C]$ are uniquely given by $M$ (not just up to null sets) for all $C \ins V$.
    So every equality between such Markov kernels, like $\Qf[\idc S] = \Qf[S]$, is unambiguous.
    The only source of ambiguity are the Markov kernels $Q(X_v|X_{\MBl^G_<(v)})$, which are only determined up to some form of null set.
    However, they are always chosen w.r.t.\ a (sub-)CADMG $G$, a topological order $<$ and one of its district $D \in \Dst[V]$ such that 
    $Q(X_v|X_{\MBl^G_<(v)})$ is a version of the marginal conditionals of both $\Qf[V]$ and $\Qf[D]$:
 \[ P(X_v|X_{\Pred^{[V]}_<(v)},\doit(X_J)) \qquad \text{and} \qquad P(X_v|X_{\Pred^{[D]}_<(v)},\doit(X_{J \cup V\sm D})).\]
 So if $<'$ is another topological order for $G$ and $Q'(X_v|X_{F(v)})$ another such version for each $v \in D$ 
 w.r.t.\ $G$, $<'$ and $D$ 
 like above, then we get the (sure) equality by the chain rule (or just the definition of conditional Markov kernels):
  \begin{align*}
      \bigotimes_{v \in D}^{>'}  Q'(X_v|X_{F(v)}) &= \Qf[D] = \bigotimes_{v \in D}^> Q(X_v|X_{\MBl^G_<(v)}).
  \end{align*}
    So, even though, the topological orders $<$ of $G$ and the $Q(X_v|X_{\MBl^G_<(v)})$'s are ambiguous up to some null sets, their product is not, when restricted to these joint choices. 
    Similarly, the final product:
    \[ \Qf[B^C] = \lB \bigotimes_{S_k \in \Dst[B^C]}^> \rB \Qf[S_k],\]
    is unambiguous, as well as the equality $\Qf[\idc S_k]=\Qf[S_k]$, and its marginal $P(X_B|\doit(X_{J \cup C}))$.

    This shows that in the non-FAIL case $P(X_B|\doit(X_{J \cup C}))$ is \trackable\ up to oracle choices from $P(X_V|\doit(X_J))$ and $G$ alone.
\end{comment}
    \end{proof}
\end{Cor}

\begin{Thm}[Soundness up to null-sets]
    \label{thm:ID-as-soundness}
    Let $G=(J,V,E,L)$ be a CADMG with a fixed topological order $<$. 
    Consider the class of IL-CBNs $M$ with observational CADMG $G$ such that the following holds: 
    \begin{enumerate}
        \item Each measurable space $\Xcal_v$ comes equipped with a fixed measure $\mu_v$ for $v \in V$,
        \item for every subset $D \ins V$ the interventional Markov kernel $\Qf[D]=P(X_D|\doit(X_{J \cup V \sm D}))$ is
            \emph{absolute continuous} w.r.t.\ the product measure $\mu_D:=\bigotimes_{v \in D} \mu_v$ and vice versa:
            \[ \mu_D \ll \Qf[D] \ll \mu_D.   \]
    \end{enumerate}
    If the ID-algorithm does not produce FAIL for input $B, C \ins G$,
    then $P(X_B|\doit(X_{J \cup C}))$ is ``almost-surely'' \trackable\ from $P(X_V|\doit(X_J))$ and $G$ for such CBNs $M$, 
    i.e.\ the Markov kernel that was output by the ID-algorithm equals $P(X_B|\doit(X_{J \cup C}))$ up to a
    $\mu_{V\sm B^C}$-null set in $\Xcal_{J \cup V\sm B^C}$ (see Remark \ref{rem:ID-as-soundness} below).
    \begin{proof}
        See Theorem \ref{thm:ID-as-soundness-prf}.
    \end{proof}
\end{Thm}

\begin{Rem}
    \label{rem:ID-as-soundness}
    \begin{enumerate}
        \item For the almost-sure soundness in Theorem \ref{thm:ID-as-soundness} to hold one implicitely needs/is allowed
            to make slight relaxations to the ID-algorithm \ref{alg:id}: 
        \item[] Instead of insisting on taking the conditionals $Q(X_v|X_{\MBl^{G'}_<(v)})$ 
            of $\Qf[D]$ for $v \in D \ins V$ one takes any version of that conditional of $\Qf[D]\otimes \mu_{V \sm D}$ 
            that is only dependent on predecessors of $v$, 
            which will always be possible as $Q(X_v|X_{\MBl^{G'}_<(v)})$ is an existing such version.
        \item The conditions of Theorem \ref{thm:ID-as-soundness} are satisfied for a IL-CBNs $M$ with CDAG
            $G^+=(J, U \dcup V,E)$ if every Markov kernel $P_v(X_v|\doit(X_{\Pa^{G^+}(v)}))$ has a strictly positive Doob-Radon-Nikodym derivative/density w.r.t.\ the measure $\mu_v$ for all $v \in V$, see Lemma \ref{lem:cbn-str-pos-density}.
            %and some additional measures $\mu_v$ for $v \in U$.
    \end{enumerate}
\end{Rem}

\begin{Thm}[From almost-sure to sure soundness]
    \label{thm:ID-as-soundness-cont}
    In addition to the conditions in Theorem \ref{thm:ID-as-soundness} assume:
          \begin{enumerate}
              \item $\Xcal_v$ is a \emph{Polish space} for every $v \in J \cup V$,
              \item $\mu_{V}$ is \emph{strictly positive} (on non-empty open subsets of $\Xcal_V$),
              %\item $\mu_{V\sm B^C}$ is \emph{strictly positive} (on non-empty open subsets),
              \item the queried interventional Markov kernel is \emph{continuous} as a map:
                  \[ P(X_B|\doit(X_{J \cup C})):\, \Xcal_{J \cup C} \to \Prob(\Xcal_B).\]
          \end{enumerate}
          Then the output $\hat P(X_B|X_{J \cup V \sm B^C})$ of the ID-algorithm (in the not-FAIL case) 
          can be changed on a $\mu_{V\sm B^C}$-null set in $\Xcal_{J \cup V \sm B^C}$ such that it becomes continuous as a map:
        \[ \hat P(X_B|X_{J \cup V \sm B^C}):\, \Xcal_{J \cup V \sm B^C} \to \Prob(\Xcal_B). \]
          Every such continuous version of $\hat P(X_B|X_{J \cup V \sm B^C})$ is then necessarily identical to
          the interventional Markov kernel $P(X_B|\doit(X_{J \cup C}))$. 
          These conditions thus allow us to recover from the ambiguity resulting from the null-sets.
          \begin{proof}
              The existence of such a null-set is clear, 
              because $\hat P(X_B|X_{J \cup V \sm B^C})$ is $\mu_{V\sm B^C}$-almost-surely equal to
              $P(X_B|\doit(X_{J \cup C}))$, and the latter was assumed to be continuous.
              The uniqueness follows from Lemma \ref{lem:cont-cond-unique}
          \end{proof}
\end{Thm}

\begin{Rem}
    The statement of Theorem \ref{thm:ID-as-soundness-cont} 
    can be further relaxed by asking for Polish spaces $\Xcal_v$ only for $v \in V$,
          strict positivity only for $\mu_{V\sm B^C}$ and only for the continuity of the maps:
          \[  \Xcal_{C \sm J} \to \Prob(\Xcal_B),\quad x_{C\sm J} \mapsto P(X_B|\doit(X_{J \cup C}=(x_J,x_{C \sm J})), \]
          for every $x_J \in \Xcal_J$ separately, by then applying the criterion from Theorem \ref{thm:ID-as-soundness-cont} 
          for each partial input $x_J\in \Xcal_J$ separately.
\end{Rem}

\begin{Eg}
    All stated assumptions of Theorems \ref{thm:ID-as-soundness} and \ref{thm:ID-as-soundness-cont} 
    are satisfied if every Markov kernel of the IL-CBN $M$ is linear Gaussian: 
    \[P_v(X_v \in dx_v|\doit(X_{\Pa^{G^+}(v)}=x_{\Pa^{G^+}(v)})) = \Ncal(dx_v|\Gamma_v \cdot x_{\Pa^{G^+}(v)} + \gamma_v, \Sigma_v), \] 
    with transition matrix $\Gamma_v$, translation vector $\gamma_v$ and positive definite covariance matrix $\Sigma_v \succ 0$ and Lebesgue measures $\mu_v$, $v \in U \cup V$.
\end{Eg}

\begin{Thm}[Completeness, see \cite{HV08s}] 
    The ID-algorithm is complete.

    More precisely, if the ID-algorithm outputs FAIL for subsets $B, C \ins G=(J,V,E,L)$ 
    then there exist two IL-CBNs $M_1$ and $M_2$ with the same observational CADMG
    $G$, the same and discrete underlying spaces $\Xcal_v$ for $v \in J \cup V$, and the same 
    observational Markov kernels $P_1(X_V|\doit(X_J)) =P_2(X_V|\doit(X_J))$ that have strictly positive mass functions such that:
    \[ P_1(X_B|\doit(X_{J \cup C})) \neq P_2(X_B|\doit(X_{J \cup C})). \]
    In particular, in case of FAIL, $P(X_B|\doit(X_{J \cup C}))$ is not identifiable from $P(X_V|\doit(X_J))$ and $G$.
\end{Thm}


\begin{Rem}[Identification of conditional causal effects, see \cite{Tian04s}] 
    If we want to know if the conditional interventional Markov kernel $P(X_A|X_B,\doit(X_{J \cup C}))$
    is \trackable\ up to oracle choices from $P(X_V|\doit(X_J))$ and $G$ then we can run the ID-algorithm
    for $A \cup B$ and $C$. If it does not output FAIL then $P(X_{A \cup B}|\doit(X_{J \cup C}))$ 
    is \trackable\ up to oracle choices from $P(X_V|\doit(X_J))$ and $G$ and by conditioning on $X_B$ afterwards so will 
    $P(X_A|X_B,\doit(X_{J \cup C}))$ be.

    However, note that there is a ``conditional'' version of the ID-algorithm, see \cite{Tian04s}, 
    that can check if (and conclude that)
    $P(X_A|X_B,\doit(X_{J \cup C}))$ is \trackable\ up to oracle choices from $P(X_V|\doit(X_J))$ and $G$ 
    even if the (unconditional) ID-algorithm outputs FAIL for $P(X_{A \cup B}|\doit(X_{J \cup C}))$.
\end{Rem}

\begin{comment}
\begin{Def}[The ID-components]
    Let $G=(J,V,E,L)$ be a CADMG. 
    \begin{enumerate}
        \item Let $B \ins V$ be a non-empty subset and $S \in \Dst[B]$ be a district. 
        \item[] Iteratively define: $A_{S,0}:=V$ and for $n \in \N$:
            \[ D_{S,n} := \Anc^{[A_{S,n}]}(S), \qquad A_{S,n+1}:=\Dist^{[D_{S,n}]}(S), \qquad  \idc S := \bigcap_{n \in \N} A_{S,n} \sni S. \]
            We then define the set of \emph{ID-components} of $B$ as:
            \[ \IDC[B]:=\lC \idc S \st S \in \Dst[B] \rC. \]
        \item Let $\emptyset \neq B \ins V$ and $C \ins J \cup V$ disjoint subsets. We then put:
            \[B^C := \Anc^{G_{\doit(C)}}(B) \sm (J \cup C) \ins V. \]
            With this we also have the corresponding sets: $\Dst[B^C]$ and $\IDC[B^C]$.
            %\begin{align*} \Dst[B|\doit(C)] &:= \Dst\lB \Anc^{G_{\doit(C)}}(B) \sm (J \cup C) \rB, \\
            %    \IDC[B|\doit(C)] &:= \IDC\lB \Anc^{G_{\doit(C)}}(B) \sm (J \cup C) \rB. 
           % \end{align*}
    \end{enumerate}
\end{Def}

\begin{Lem}[ID-components are \trackable\ up to oracle choices]
    Let $M$ be an IL-CBN with with observational CADMG $G=(J,V,E,L)$ and a fixed topological order $<$.
    Let $B \ins V$ any (non-empty) subset and $\idc S \in \IDC[B]$ any ID-component of $B$.
    Then the interventional distribution $\Qf[\idc S]$ is \trackable\ from $\Qf[V]$ up to oracle choices.
    \begin{proof}
        This follows recursively from Remark \ref{rem:Q-dist-ident} and Lemma \ref{lem:int-anc}.
    \end{proof}
\end{Lem}

\begin{Cor}[ID-algorithm rephrased]
    Let $M$ be an IL-CBN with with observational CADMG $G=(J,V,E,L)$ and a fixed topological order $<$.
    Let $\emptyset \neq B \ins V$ and $C \ins J \cup V$ be two disjoint subsets of nodes. 
    Then the interventional Markov kernel $P(X_B|\doit(X_{J \cup C}))$ is \trackable\ up to oracle choices from 
    the observational Markov kernel $P(X_V|\doit(X_J))$ and $G$ if we have the inclusion:
    \[\Dst[B^C]  \ins \IDC[B^C], \]
    or, in other words, if $S=\idc S$ for every $S \in \Dst[B^C]$.
    If this is the case then $P(X_B|\doit(X_{J \cup C}))$ is the $B$-marginal of:
    %\[ \lB \bigotimes_{\idc S \in \IDC[B^C] }^> \rB \Qf[\idc S].  \]
    \[ \Qf[B^C] = \lB \bigotimes_{S \in \Dst[B^C] }^> \rB \Qf[S] = \lB \bigotimes_{\idc S \in \IDC[B^C] }^> \rB \Qf[\idc S].  \]
\end{Cor}

\end{comment}


\paragraph{Examples}
%\subsubsection{Examples}


\begin{Eg}
    Consider the DAG $G=(V,E)$ from Figure \ref{fig:dag-2} (a) with $V=\lC v_1,v_2 \rC$, 
    $E=\lC v_1 \tuh v_2 \rC$.
     We want to determine if we can identify $P(X_2|\doit(X_1))$ from $P(X_1,X_2)$ in case we have a discrete CBN with strictly positive mass function: $p(x_1,x_2,x_3) >0$.  
     For this let $B:=\lC v_2 \rC$ and $C:=\lC v_1 \rC$. Note that we have the topological order $v_1 < v_2$.
    We then follow the steps of the ID-algorithm:
    \begin{enumerate}
        \item $\Qf[V](x_1,x_2) := p(x_1,x_2)$.
        \item $B^C = \Anc^{G_{\doit(C)}}(B) \sm C = \lC v_2\rC$.
        \item $\Dst[B^C] =\lC S=\lC v_2 \rC \rC$. So we compute:
            \begin{enumerate}
                \item $D_{0} = \Dist^{[V]}(S) = \lC v_2 \rC$. Compute:
                    \begin{align*}
                        \Qf[D_{0}](x_2|x_1) &= q(x_2|x_1) = \frac{\Qf[V](x_1,x_2)}{\Qf[V](x_1)} = p(x_2|x_1).
                    \end{align*}
                \item $A_{1} = \Anc^{[D_{0}]}(S) = \lC v_2\rC=D_{0}$, thus $\idc S = D_{0}=\lC v_2 \rC =S$. Compute:
                    \begin{align*}
                        \Qf[S](x_2|x_1) = \Qf[D_{0}](x_2|x_1) = p(x_2|x_1). 
                    \end{align*}
            \end{enumerate}
        \item Since $\idc S = S = \lC v_2 \rC$ we can compute: 
            \begin{align*}
                   \Qf[B^C](x_2|x_1)  = \Qf[S_1](x_2|x_1) = p(x_2|x_1) . 
               \end{align*}
    \end{enumerate}
    So we can identify $P(X_2|\doit(X_1))$ from $P(X_1,X_2)$ as the conditional $P(X_2|X_1)$ via the mass function from above.
\end{Eg}

\begin{Eg}
    Consider the ADMG $G=(V,E,L)$ from Figure \ref{fig:dag-2} (b) with $V=\lC v_1,v_2 \rC$, 
    $E=\lC v_1 \tuh v_2 \rC$ and $L=\lC v_1 \huh v_2 \rC$.
     We want to determine if we can identify $P(X_2|\doit(X_1))$ from $P(X_1,X_2)$ in case we have a discrete CBN with strictly positive mass function: $p(x_1,x_2,x_3) >0$.  
     For this let $B:=\lC v_2 \rC$ and $C:=\lC v_1 \rC$. Note that we have the topological order $v_1 < v_2$.
    We then follow the steps of the ID-algorithm:
    \begin{enumerate}
        \item $\Qf[V](x_1,x_2) := p(x_1,x_2)$.
        \item $B^C = \Anc^{G_{\doit(C)}}(B) \sm C = \lC v_2\rC$.
        \item $\Dst[B^C] =\lC S=\lC v_2 \rC \rC$. So we compute:
            \begin{enumerate}
                \item $D_{0} = \Dist^{[V]}(S) = \lC v_1,v_2 \rC$. Compute:
                    \begin{align*}
                        q(x_1) &=  \Qf[V](x_1) &&= p(x_1), \\
                        q(x_2|x_1) &= \frac{\Qf[V](x_1,x_2)}{\Qf[V](x_1)} && = p(x_2|x_1), \\
                        \Qf[D_{0}](x_1,x_2) &= q(x_2|x_1) \cdot q(x_1) &&= p(x_1,x_2).
                    \end{align*}
                \item $A_{1} = \Anc^{[D_{0}]}(S) = \lC v_1,v_2\rC=D_{0}$, thus $\idc S = D_{0}=\lC v_1,v_2 \rC$. Compute:
                    \begin{align*}
                        \Qf[\idc S](x_1,x_2) = \Qf[D_{0}](x_1,x_2) = p(x_1,x_2). 
                    \end{align*}
            \end{enumerate}
        \item Since $\idc S = \lC v_1,v_2 \rC \neq \lC v_2 \rC = S$ the ID-algorithms outputs: FAIL.
    \end{enumerate}
    So we can not identify $P(X_2|\doit(X_1))$ from $P(X_1,X_2)$ and $G$.
\end{Eg}

\begin{figure}[ht]
\centering
\begin{tikzpicture}[scale=1, transform shape]
   \begin{scope}
    \node[ndout] (v1) at (0,0) {$v_1$};
    \node[ndout] (v2) at (-1.5,-2) {$v_2$};
    \node[ndout] (v3) at (1.5,-2) {$v_3$};
    \draw[arout] (v1) to (v2);
    \draw[arout] (v1) to (v3);
    \draw[arout] (v2) to (v3);
    \node (a) at (-1.7,0.2) {(a)};
  \end{scope} 
   \begin{scope}[xshift=5.5cm]
    \node[ndout] (v1) at (0,0) {$v_1$};
    \node[ndint] (v2) at (-1.5,-2) {$v_2$};
    \node[ndout] (v3) at (1.5,-2) {$v_3$};
    \draw[arout] (v1) to (v3);
    \draw[arint] (v2) to (v3);
    \node (a) at (-1.7,0.2) {(b)};
  \end{scope} 
  \begin{scope}[xshift=11cm]
    \node[ndint] (v1) at (0,0) {$v_1$};
    \node[ndint] (v2) at (-1.5,-2) {$v_2$};
    \node[ndout] (v3) at (1.5,-2) {$v_3$};
    \draw[arint] (v1) to (v3);
    \draw[arint] (v2) to (v3);
    \node (a) at (-1.7,0.2) {(c)};
  \end{scope}
\end{tikzpicture}
\caption{ A DAG with three nodes and its intervened graphs.}
\label{fig:dag-3}
\end{figure}

\begin{Eg}
   Consider the DAG $G=(V,E,L)$ from Figure \ref{fig:dag-3} with $V=\lC v_1,v_2,v_3 \rC$, 
    $E=\lC v_1 \tuh v_2, v_1 \tuh v_3, v_2 \tuh v_3  \rC$.
     We want to determine if we can identify $P(X_3|\doit(X_2))$ from $P(X_1,X_2,X_3)$ in case we have a discrete CBN with strictly positive mass function: $p(x_1,x_2,x_3) >0$.  
     For this let $B:=\lC v_3 \rC$ and $C:=\lC v_2 \rC$. Note that we have the topological order $v_1 < v_2 < v_3$.
    We then follow the steps of the ID-algorithm:
        \begin{enumerate}
        \item $\Qf[V](x_1,x_2,x_3) := p(x_1,x_2,x_3)$.
        \item $B^C = \Anc^{G_{\doit(C)}}(B) \sm C = \lC v_1,v_3 \rC$.
        \item $\Dst[B^C] =\lC S_1=\lC v_1 \rC, S_2=\lC v_3 \rC   \rC$, see Figure \ref{fig:dag-3} (b).
        \item For $S_1=\lC v_1 \rC$:
            \begin{enumerate}
                \item $D_{1,0} = \Dist^{[V]}(S_1) = \lC v_1\rC$. Compute:
                    \begin{align*}
                        \Qf[D_{1,0}](x_1) &= q(x_1) = \Qf[V](x_1)= p(x_1).
                    \end{align*}
                \item $A_{1,1} = \Anc^{[D_{1,0}]}(S_1) = \lC v_1\rC = D_{1,0}$, thus $\idc S_1 =D_{1,0}$. Compute:
                    \begin{align*}
                        \Qf[\idc S_1](x_1) =    \Qf[D_{1,0}](x_1)   = p(x_1).
                    \end{align*}
                \item $ S_1 = \lC v_1 \rC = \idc S_1$. Compute:
                    \begin{align*}
                       \Qf[S_1](x_1) =    \Qf[\idc S_1](x_1) =  p(x_1).
                    \end{align*}
            \end{enumerate}
       \item For $S_2=\lC v_3 \rC$:
            \begin{enumerate}
                \item $D_{2,0} = \Dist^{[V]}(S_2) = \lC v_3 \rC$. Compute:
                    \begin{align*}
                \Qf[D_{2,0}](x_3|x_1,x_2) =  q(x_3|x_1,x_2)  = \frac{\Qf[V](x_1,x_2,x_3)}{\Qf[V](x_1,x_2)}  = p(x_3|x_1,x_2).
                    \end{align*}
                \item $A_{2,1} = \Anc^{[D_{2,0}]}(S_1) = \lC v_3  \rC = D_{2,0}$. thus $\idc S_2 = D_{2,0}$. Compute:
                 \begin{align*}
               \Qf[\idc S_2](x_3|x_1,x_2)  = \Qf[D_{2,0}](x_3|x_1,x_2)  = p(x_3|x_1,x_2).
                    \end{align*}
                \item $S_2 = \lC v_3 \rC =\idc S_2$. Compute:
                 \begin{align*}
               \Qf[S_2](x_3|x_1,x_2)  = \Qf[\idc S_2](x_3|x_1,x_2)  = p(x_3|x_1,x_2).
                    \end{align*}
            \end{enumerate}
        \item Since both $\idc S_1 = S_1 = \lC v_1 \rC$ and $\idc S_2 =S_2 = \lC v_3\rC$ we can compute:
               \begin{align*}
                   \Qf[B^C](x_1,x_3|x_2) &= \Qf[S_2](x_3|x_1,x_2) \cdot \Qf[S_1](x_1) \\ 
                                         &= p(x_3|x_1,x_2) \cdot p(x_1), \\
                   p(x_3|\doit(x_2)) & = \sum_{x_1} \Qf[B^C](x_1,x_3|x_2) \\ 
                   & = \sum_{x_1} p(x_3|x_1,x_2) \cdot p(x_1). 
               \end{align*}
    \end{enumerate}
    So we can identify $P(X_3|\doit(X_2))$ from $P(X_1,X_2,X_3)$ via the mass function from above.
\end{Eg}

\begin{Eg}[Counter example when mass functions are not strictly positive]\label{eg:counterexample_nonpositive_backdoor_adjustment}
   Consider the DAG $G=(V,E,L)$ from Figure \ref{fig:dag-3} with $V=\lC v_1,v_2,v_3 \rC$, 
    $E=\lC v_1 \tuh v_2, v_1 \tuh v_3, v_2 \tuh v_3  \rC$.
     We want to determine if we can identify $P(X_3|\doit(X_2))$ from $P(X_1,X_2,X_3)$ in case we do NOT have a strictly positive mass function: $p(x_1,x_2,x_3) >0$. We assume $\Xcal_1=\Xcal_2=\Xcal_3 = \lC 0,1 \rC$.
     \begin{align*}
         p(x_1=1) & := \frac{1}{2}, &
         p(x_2|\doit(x_1)) & := \delta_{x_1}(x_2),\\
         p(x_3=1|\doit(x_1=0,x_2=0)) &:= \frac{1}{4}, & 
         p(x_3=1|\doit(x_1=1,x_2=0)) &:= \frac{3}{4}, \\
         p(x_3=1|\doit(x_1=0,x_2=1)) &:= \frac{1}{8}, &
         p(x_3=1|\doit(x_1=1,x_2=1)) &:= \frac{3}{8}.
     \end{align*}
     Then we get for the $(X_1,X_2)$-marginal of the observational distribution
     $P(X_1,X_2) = P(X_2|\doit(X_1)) \otimes P(X_1)$ the mass functions:
    \begin{align*}
         p(x_1=0,x_2=0) &=  p(x_1=1,x_2=1) = \frac{1}{2}, \\
         p(x_1=1,x_2=0) &=  p(x_1=0,x_2=1) = 0.
    \end{align*}
    Note that this shows that $P(X_1,X_2)$ and thus $P(X_1,X_2,X_3)$ do not have a strictly positive mass functions.
    With this a valid conditional for the observational joint distribution $P(X_1,X_2,X_3)$ conditioned on $(X_1,X_2)$ is:
     \begin{align*}
         p(x_3=1|x_1=0,x_2=0) &:= \frac{1}{4}, & 
         p(x_3=1|x_1=1,x_2=0) &:= \frac{3}{8}, \\
         p(x_3=1|x_1=0,x_2=1) &:= \frac{5}{8}, &
         p(x_3=1|x_1=1,x_2=1) &:= \frac{3}{8}.
     \end{align*}
    The interventional distribution is given by:
     \begin{align*}
         p(x_3|\doit(x_2)) & = \sum_{x_1} p(x_3|\doit(x_1,x_2)) \cdot p(x_1),\\
          %  & = \frac{1}{2} \lp p(x_3|\doit(x_1=0,x_2)) + p(x_3|\doit(x_1=1,x_2)) \rp,  \\ 
            p(x_3=1|\doit(x_2=0)) & = \frac{1}{2} \lp p(x_3=1|\doit(x_1=0,x_2=0)) + p(x_3=1|\doit(x_1=1,x_2=0)) \rp  \\
             &=\frac{1}{2} \lp \frac{1}{4} + \frac{3}{4}  \rp = \frac{1}{2} ,\\
            p(x_3=1|\doit(x_2=1)) & = \frac{1}{2} \lp p(x_3=1|\doit(x_1=0,x_2=1)) + p(x_3=1|\doit(x_1=1,x_2=1)) \rp  \\ 
            &= \frac{1}{2} \lp \frac{1}{8} + \frac{3}{8}  \rp = \frac{1}{4}.
    \end{align*}
    On the other hand, using the other conditional mass functions instead, gives us:
    \begin{align*}
         \hat p(x_3|x_2) & := \sum_{x_1} p(x_3|x_1,x_2) \cdot p(x_1),\\
         \hat p(x_3=1|x_2=0) & = \frac{1}{2} \lp p(x_3=1|x_1=0,x_2=0) + p(x_3=1|x_1=1,x_2=0) \rp  \\
         &= \frac{1}{2} \lp \frac{1}{4} + \frac{1}{2} \rp = \frac{3}{8} \\& \neq \frac{1}{2} =p(x_3=1|\doit(x_2=0)), \\
         \hat p(x_3=1|x_2=1) & = \frac{1}{2} \lp p(x_3=1|x_1=0,x_2=1) + p(x_3=1|x_1=1,x_2=1) \rp  \\ 
         &= \frac{1}{2} \lp \frac{5}{8} + \frac{1}{8}  \rp = \frac{3}{8} \\& \neq \frac{1}{4} =p(x_3=1|\doit(x_2=1)).
     \end{align*}
     This shows that the ``surrogate'' Markov kernel $\hat P(X_3|X_2) := P(X_3|X_1,X_2) \circ P(X_1)$, 
     which would be proposed by both the ID-algorithm and the backdoor criterion, 
     is NOT equal to the interventional Markov kernel
     $P(X_3|\doit(X_2)) = P(X_3|\doit(X_1,X_2)) \circ P(X_1) $, not even $P(X_2)$-almost-surely.
 \end{Eg}

\begin{figure}[ht]
\centering
\begin{tikzpicture}[scale=1, transform shape]
  \begin{scope}
    \node[ndout] (v1) at (0,0) {$v_1$};
    \node[ndout] (v2) at (2,0) {$v_2$};
    \node[ndout] (v3) at (4,0) {$v_3$};
    \draw[arout] (v1) to (v2);
    \draw[arout] (v2) to (v3);
    \draw[arlat, bend left] (v1) to (v3);
    \node (a) at (-1,0.5) {(a)};
  \end{scope}
 \begin{scope}[xshift=7cm]
    \node[ndint] (v1) at (0,0) {$v_1$};
    \node[ndout] (v2) at (2,0) {$v_2$};
    \node[ndout] (v3) at (4,0) {$v_3$};
    \draw[arint] (v1) to (v2);
    \draw[arout] (v2) to (v3);
    \node (a) at (-1,0.5) {(b)};
  \end{scope}
  \begin{scope}[yshift=-2cm]
    \node[ndint] (v1) at (0,0) {$v_1$};
    \node[ndint] (v2) at (2,0) {$v_2$};
    \node[ndout] (v3) at (4,0) {$v_3$};
    \draw[arint] (v2) to (v3);
    \node (a) at (-1,0.5) {(c)};
  \end{scope}
  \begin{scope}[yshift=-2cm,xshift=7cm]
    \node[ndint] (v1) at (0,0) {$v_1$};
    \node[ndout] (v2) at (2,0) {$v_2$};
    \node[ndint] (v3) at (4,0) {$v_3$};
    \draw[arint] (v1) to (v2);
    \node (a) at (-1,0.5) {(d)};
  \end{scope}
  \begin{scope}[yshift=-4cm]
    \node[ndout] (v1) at (0,0) {$v_1$};
    \node[ndout] (v2) at (2,0) {$v_2$};
    \node[ndint] (v3) at (4,0) {$v_3$};
    \draw[arout] (v1) to (v2);
    \node (a) at (-1,0.5) {(e)};
  \end{scope}
  \begin{scope}[yshift=-4cm,xshift=7cm]
    \node[ndout] (v1) at (0,0) {$v_1$};
    \node[ndint] (v2) at (2,0) {$v_2$};
    \node[ndout] (v3) at (4,0) {$v_3$};
    \draw[arint] (v2) to (v3);
    \draw[arlat, bend left] (v1) to (v3);
    \node (a) at (-1,0.5) {(f)};
  \end{scope}
\end{tikzpicture}
\caption{An ADMG and its mutilations, corresponding to the interventional Markov kernels: (a) $\Qf[\lC v_1,v_2,v_3 \rC]$, (b) $\Qf[\lC v_2, v_3 \rC]$, (c) $\Qf[\lC v_3 \rC]$, (d) $\Qf[\lC v_2 \rC]$, (e) $\Qf[\lC v_1, v_2 \rC]$, (f)  $\Qf[\lC v_2, v_3 \rC]$. }
\label{fig:admg-2}
\end{figure}



\begin{Eg}
    Consider the ADMG $G=(V,E,L)$ from Figure \ref{fig:admg-2} with $V=\lC v_1,v_2,v_3 \rC$, 
    $E=\lC v_1 \tuh v_2, v_2 \tuh v_3 \rC$, $J=\emptyset$ and $L=\lC v_1 \huh v_3 \rC$.
     We want to determine if we can identify $P(X_3|\doit(X_1))$ from $P(X_1,X_2,X_3)$ in case we have a discrete CBN with strictly positive mass function: $p(x_1,x_2,x_3) >0$.  
     For this let $B:=\lC v_3 \rC$ and $C:=\lC v_1 \rC$. Note that we have the topological order $v_1 < v_2 < v_3$.
    We then follow the steps of the ID-algorithm:
    \begin{enumerate}
        \item $\Qf[V](x_1,x_2,x_3) := p(x_1,x_2,x_3)$.
        \item $B^C = \Anc^{G_{\doit(C)}}(B) \sm C = \lC v_2,v_3 \rC$.
        \item $\Dst[B^C] =\lC S_1=\lC v_3 \rC, S_2=\lC v_2 \rC   \rC$, see Figure \ref{fig:admg-2} (b).
        \item For $S_1=\lC v_3 \rC$:
            \begin{enumerate}
                \item $D_{1,0} = \Dist^{[V]}(S_1) = \lC v_1, v_3\rC$, see Figure \ref{fig:admg-2} (a), (f). Compute:
                    \begin{align*}
                        q(x_1) & = \Qf[V](x_1) &&= p(x_1),\\
                        q(x_3|x_1,x_2) & = \frac{\Qf[V](x_1,x_2,x_3)}{\Qf[V](x_1,x_2)} && = p(x_3|x_1,x_2), \\
                        \Qf[D_{1,0}](x_1,x_3|x_2) &= q(x_3|x_1,x_2) \cdot q(x_1) && = p(x_3|x_1,x_2) \cdot p(x_1).
                    \end{align*}
                \item $A_{1,1} = \Anc^{[D_{1,0}]}(S_1) = \lC v_3\rC$. see Figure \ref{fig:admg-2} (f), (c). Compute:
                    \begin{align*}
                        \Qf[A_{1,1}](x_3|x_2) = \sum_{x_1} \Qf[D_{1,0}](x_1,x_3|x_2) = \sum_{x_1} p(x_3|x_1,x_2) \cdot p(x_1).
                    \end{align*}
                \item $D_{1,1} = \Dist^{[A_{1,1}]}(S_1) = \lC v_3\rC = A_{1,1}$, thus $\idc S_1 =A_{1,1}$ . Compute:
                   \begin{align*}
                       \Qf[\idc S_1](x_3|x_2) =    \Qf[A_{1,1}](x_3|x_2) = \sum_{x_1} p(x_3|x_1,x_2) \cdot p(x_1).
                    \end{align*}
                \item $ S_1 = \lC v_3 \rC = \idc S_1$. Compute:
                    \begin{align*}
                       \Qf[S_1](x_3|x_2) =    \Qf[\idc S_1](x_3|x_2) = \sum_{x_1} p(x_3|x_1,x_2) \cdot p(x_1).
                    \end{align*}
            \end{enumerate}
       \item For $S_2=\lC v_2 \rC$:
            \begin{enumerate}
                \item $D_{2,0} = \Dist^{[V]}(S_2) = \lC v_2 \rC$, see Figure \ref{fig:admg-2} (a), (d). Compute:
                    \begin{align*}
                \Qf[D_{2,0}](x_2|x_1) =  q(x_2|x_1)  = \frac{\Qf[V](x_1,x_2)}{\Qf[V](x_1)}  = p(x_2|x_1).
                    \end{align*}
                \item $A_{2,1} = \Anc^{[D_{2,0}]}(S_1) = \lC v_2  \rC = D_{2,0}$. thus $\idc S_2 = D_{2,0}$. Compute:
                 \begin{align*}
               \Qf[\idc S_2](x_2|x_1)  = \Qf[D_{2,0}](x_2|x_1)  = p(x_2|x_1).
                    \end{align*}
                \item $S_2 = \lC v_2 \rC =\idc S_2$. Compute:
                 \begin{align*}
               \Qf[S_2](x_2|x_1)  = \Qf[\idc S_2](x_2|x_1)  = p(x_2|x_1).
                    \end{align*}
            \end{enumerate}
        \item Since both $\idc S_1 = S_1 = \lC v_3 \rC$ and $\idc S_2 =S_2 = \lC v_2\rC$ we can compute:
               \begin{align*}
                   \Qf[B^C](x_2,x_3|x_1) &= \Qf[S_2](x_2|x_1) \cdot \Qf[S_1](x_3|x_2) \\ 
                                         &= p(x_2|x_1) \cdot \sum_{x_1'} p(x_3|x_1',x_2) \cdot p(x_1'), \\
                   p(x_3|\doit(x_1)) & = \sum_{x_2} \Qf[B^C](x_2,x_3|x_1) \\ 
                   & = \sum_{x_2} p(x_2|x_1) \cdot \sum_{x_1'} p(x_3|x_1',x_2) \cdot p(x_1') \\
                   & = \sum_{x_1',x_2'} p(x_3|x_1',x_2') \cdot p(x_1') \cdot p(x_2'|x_1).   
               \end{align*}
    \end{enumerate}
    So we can identify $P(X_3|\doit(X_1))$ from $P(X_1,X_2,X_3)$ via the mass function from above.
\end{Eg}









\paragraph{Proofs - Soundness Criteria}
%\subsubsection{Soundness Criteria}


We have seen in Corollary \ref{cor:id-alg-sound-oracle} that the \emph{ID-Algorithm} \ref{alg:id} is \emph{sound up to oracle choices}.
In this subsection we want to investigate the possibility of other forms of soundness that would allow for stronger forms of identifiability and/or \trackability.






\begin{Thm}[Soundness up to null-sets]
    \label{thm:ID-as-soundness-prf}
    Let $G=(J,V,E,L)$ be a CADMG with a fixed topological order $<$. 
    Consider the class of IL-CBNs $M$ with observational CADMG $G$ such that the following holds: 
    \begin{enumerate}
        \item The measurable spaces $\Xcal_v$ come equipped with a measure $\mu_v$, $v \in V$,
        \item for every subset $D \ins V$ the interventional Markov kernel $\Qf[D]=P(X_D|\doit(X_{J \cup V \sm D}))$ is
            \emph{absolute continuous} w.r.t.\ the product measure $\mu_D:=\otimes_{v \in D} \mu_v$ and vice versa:
            \[ \mu_D \ll \Qf[D] \ll \mu_D.   \]
    \end{enumerate}
    If the ID-algorithm does not produce FAIL for input $B, C \ins G$,
    then $P(X_B|\doit(X_{J \cup C}))$ is ``almost-surely'' \trackable\ from $P(X_V|\doit(X_J))$ and $G$ for such CBNs $M$, 
    i.e.\ the Markov kernel that was output by the ID-algorithm equals $P(X_B|\doit(X_{J \cup C}))$ up to a
    $\mu_{V\sm B^C}$-null set in $\Xcal_{J \cup V\sm B^C}$.
    \begin{proof}
        Since by the second assumption we have for each $\mu_v$ and any fixed value $x_{\lC v \rC^\cmpl}$:
        \[  P(X_v|\doit(X_{\lC v \rC^\cmpl}=x_{\lC v \rC^\cmpl})) \ll \mu_v \ll P(X_v|\doit(X_{\lC v \rC^\cmpl}=x_{\lC v \rC^\cmpl})),  \]
        we can w.l.o.g.\ assume that the $\mu_v$ are probability measures for $v \in V$.
       
        a) Now consider any subset $A \ins V$ and $D \in \Dst[A]$ and $v \in D$.  We abbreviate:
        \[ G':= G_{\doit(V \sm A)}, \qquad A_<:=\Pred^{[A]}_<(v), \qquad D_<:=\Pred^{[D]}_<(v). \]
         Assume that we have a Markov kernel:
        \[ \Kcal[A]=\tilde P(X_{A}|\doit(X_{J \cup V \sm A})) :\, \Xcal_{J \cup V \sm A} \dshto \Xcal_A,  \]
        such that:
        \[ \Kcal[A] = \Qf[A]\qquad \mu_{V\sm A}\text{-a.s.}  \]
         Note that the almost sure equality from above implies the equality:
        \begin{align}
            \Kcal[A] \otimes \mu_{V\sm A} = \Qf[A]\otimes \mu_{V\sm A}. \label{eq:K-Q-as} 
        \end{align}
        We want to show that if we perfom the steps of the ID-algorithm that computes $\Qf[D]$ from $\Qf[A]$
        on $\Kcal[A]$ then the corresponding output, abbreviated as $\Kcal[D]$, satisfies:
        \[ \Kcal[D] = \Qf[D]\qquad \mu_{V\sm D}\text{-a.s.}  \]
        For this consider any version of the conditional of the following marginal of $\Kcal[A]$:
        \[ \tilde P(X_{A_\le}|\doit(X_{J \cup V \sm A}))  \qquad \text{w.r.t.} \qquad \tilde P(X_{A_<}|\doit(X_{J \cup V \sm A})), \] 
        which we will denote by:
        \[ K(X_v|X_{J \cup V \sm A_\ge}). \]
        This by definition will satisfy:
        \begin{align}
        \tilde P(X_{A_\le}|\doit(X_{J \cup V \sm A})) =  K(X_v|X_{J \cup V \sm A_\ge}) \otimes \tilde P(X_{A_<}|\doit(X_{J \cup V \sm A})).  \label{eq:K-cond}  \end{align}
        This then implies:
        \begin{align*}
           & P(X_{A_\le}|\doit(X_{J \cup V \sm A})) \otimes \mu_{V \sm A}(X_{V \sm A}) \\
           & \stackrel{\text{Eq. }\ref{eq:K-Q-as}}{=} \tilde P(X_{A_\le}|\doit(X_{J \cup V \sm A})) \otimes \mu_{V \sm A}(X_{V \sm A}) \\
           & \stackrel{\text{Eq. }\ref{eq:K-cond}}{=} K(X_v|X_{J \cup V \sm A_\ge}) \otimes \tilde P(X_{A_<}|\doit(X_{J \cup V \sm A})) \otimes \mu_{V \sm A}(X_{V \sm A}) \\
           & \stackrel{\text{Eq. }\ref{eq:K-Q-as}}{=} K(X_v|X_{J \cup V \sm A_\ge}) \otimes  P(X_{A_<}|\doit(X_{J \cup V \sm A})) \otimes \mu_{V \sm A}(X_{V \sm A}).
        \end{align*}
        This shows that $K(X_v|X_{J \cup V \sm A_\ge})$ is a version of the conditional of:
        \[  P(X_{A_\le}|\doit(X_{J \cup V \sm A})) \otimes \mu_{V \sm A}(X_{V \sm A})  \qquad \text{w.r.t.} \qquad  P(X_{A_<}|\doit(X_{J \cup V \sm A})) \otimes \mu_{V \sm A}(X_{V \sm A}). \]
        Note that by the interventional ordered local Markov property, Proposition \ref{prp:ioLMP},
        there exists a Markov kernel $Q(X_v|X_{\MBl^{G'}_<(v)})$ that simultaneuously is a version of both:
        \[ P(X_v|X_{A_<},\doit(X_{J\cup V \sm A})) \qquad \text{and} \qquad P(X_v|X_{D_<},\doit(X_{J\cup V \sm D})), \]
        and is thus, in particular, another version of the conditional of:
        \[  P(X_{A_\le}|\doit(X_{J \cup V \sm A})) \otimes \mu_{V \sm A}(X_{V \sm A})  \qquad \text{w.r.t.} \qquad  P(X_{A_<}|\doit(X_{J \cup V \sm A})) \otimes \mu_{V \sm A}(X_{V \sm A}). \]
        Now consider the (measurable) set where these two conditional Markov kernels deviate:
        \[ \tilde N := \lC x_{J \cup V \sm A_\ge} \in \Xcal_{J \cup V \sm A_\ge}  
            \st K(X_v|X_{J \cup V \sm A_\ge}=x_{J \cup V \sm A_\ge})
        \neq Q(X_v|X_{\MBl^{G'}_<(v)} = x_{\MBl^{G'}_<(v)}) \rC, \]
        and $N := \tilde N \times \Xcal_{A_\ge} \ins \Xcal_{J \cup V}$.
        Since conditional Markov kernels are essentially unique we get that for every $x_J \in \Xcal_J$ we have:
        \[ \lp \Qf[A] \otimes \mu_{V \sm A} \rp (N^{x_J}|x_J) = 0.   \]
        Since, by assumption, we have: $\mu_A \ll \Qf[A]$, we get for every $x_J \in \Xcal_J$:
        \[ \lp \mu_D \otimes \mu_{V \sm D} \rp (N^{x_J}) = \mu_V (N^{x_J}) = \lp \mu_A \otimes \mu_{V \sm A} \rp (N^{x_J}) = 0.  \]
        Since, by assumption, we also have: $\Qf[D] \ll \mu_D$, we get for every $x_J \in \Xcal_J$:
        \[ \lp \Qf[D] \otimes \mu_{V \sm D} \rp (N^{x_J}|x_J) = 0.   \]
        %This holds in particular for $x_J':=x_J$.
        Let $\hat N := \tilde N \times \Xcal_{A_\ge \sm D_\ge}$.
        Since $D_\ge \ins A_\ge$ the set $N$ is of the form:
         \[ N = \tilde N \times \Xcal_{A_\ge} = \hat N \times \Xcal_{D_\ge}. \]
        So the above shows that we have for every $x_J \in \Xcal_J$:
        \[ \lp P(X_{D_<}|\doit(X_{J \cup V \sm D})) \otimes \mu_{V\sm D}(X_{V \sm D}) \rp (\hat N^{x_J}|x_J) = \lp \Qf[D] \otimes \mu_{V \sm D} \rp (N^{x_J}|x_J) = 0.  \]
        This shows that $K(X_v|X_{J \cup V \sm A_\ge})$ and $Q(X_v|X_{\MBl^{G'}_<(v)})$
        agree up to a measurable \\ $\lp P(X_{D_<}|\doit(X_{J \cup V \sm D})) \otimes \mu_{V\sm D}(X_{V \sm D})\rp$-null set.
        Remember that $Q(X_v|X_{\MBl^{G'}_<(v)})$ satisfies:
        \[ P(X_{D_\le}|\doit(X_{J \cup V \sm D})) = Q(X_v|X_{\MBl^{G'}_<(v)}) \otimes P(X_{D_<}|\doit(X_{J \cup V \sm D})).\]
        Together with the above we thus get:
        \begin{align*} 
            & P(X_{D_\le}|\doit(X_{J \cup V \sm D})) \otimes \mu_{V\sm D}(X_{V \sm D}) \\
        & = Q(X_v|X_{\MBl^{G'}_<(v)}) \otimes P(X_{D_<}|\doit(X_{J \cup V \sm D})) \otimes \mu_{V\sm D}(X_{V \sm D}) \\
        & = K(X_v|X_{J \cup V \sm A_\ge}) \otimes P(X_{D_<}|\doit(X_{J \cup V \sm D})) \otimes \mu_{V\sm D}(X_{V \sm D}).
    \end{align*}
       This shows that $K(X_v|X_{J \cup V \sm A_\ge})$ is version of the conditional of:
               \[ P(X_{D_\le}|\doit(X_{J \cup V \sm D})) \otimes \mu_{V\sm D}(X_{V\sm D}) \qquad \text{w.r.t.} \qquad
        P(X_{D_<}|\doit(X_{J \cup V \sm D})) \otimes \mu_{V\sm D}(X_{V\sm D}). \] 
       If we let $v$ run through $D = \lC v_1,\dots,v_K \rC$, $v_1 < v_2 < \cdots < v_K$ 
       in reverse topological order we inductively get:
       \begin{align*}
           %&  \lp \bigotimes_{v \in D}^> Q(X_v|X_{\MBl^{G'}_<(v)}) \rp \otimes \mu_{V\sm D}(X_{V\sm D})
           & \Qf[D] \otimes \mu_{V\sm D} \\
           &= P(X_D|\doit(X_{J \cup V \sm D})) \otimes \mu_{V\sm D}(X_{V\sm D}) \\
           &= K(X_{v_K}|X_{J \cup V \sm A_{\ge v_K}}) \otimes P(X_{D_{< v_{K}}}|\doit(X_{J \cup V \sm D})) \otimes \mu_{V\sm D}(X_{V\sm D})\\ 
           &= K(X_{v_K}|X_{J \cup V \sm A_{\ge v_K}}) \otimes P(X_{D_{\le v_{K-1}}}|\doit(X_{J \cup V \sm D})) \otimes \mu_{V\sm D}(X_{V\sm D})\\ 
           &= K(X_{v_K}|X_{J \cup V \sm A_{\ge v_K}}) \otimes  K(X_{v_{K-1}}|X_{J \cup V \sm A_{\ge v_{K-1}}}) \otimes P(X_{D_{< v_{K-1}}}|\doit(X_{J \cup V \sm D})) \otimes \mu_{V\sm D}(X_{V\sm D})\\ 
           & = \cdots\\
           &= \lp \bigotimes_{v \in D}^> K(X_v|X_{J \cup V \sm A_{\ge v}}) \rp \otimes \mu_{V\sm D}(X_{V\sm D}). 
       \end{align*}
       Since such factorizations are essentially unique we get that:
       \[ \Qf[D] = \bigotimes_{v \in D}^> K(X_v|X_{J \cup V \sm A_{\ge v}}) =: \Kcal[D] \qquad \mu_{V\sm D}\text{-a.s.}  \]
       %where $\mu_{V\sm D}(X_{V \sm D}|X_J)$, to be precise, indicates the constant Markov kernel:
       %\[ \mu_{V\sm D}(X_{V \sm D}|X_J):\, \Xcal_J \to \Prob(\Xcal_{V\sm D}), \qquad x_J \mapsto \mu_{V\sm D}. \]
       This shows the claim.

       b) We now reverse the situation. For a subset $A \ins V$ and every $D \in \Dst[A]=\lC D_1,\dots,D_L\rC$, 
       consider that we are given a Markov kernel:
       \[ \Kcal[D]=\tilde P(X_{D}|\doit(X_{J \cup V \sm D})) :\, \Xcal_{J \cup V \sm D} \dshto \Xcal_{D}, \]
        such that:
       \[ \Kcal[D] = \Qf[D]\qquad \mu_{V\sm D}\text{-a.s.,}  \]
       which implies the equality:
        \begin{align}
            \Kcal[D] \otimes \mu_{V\sm D} = \Qf[D]\otimes \mu_{V\sm D}. \label{eq:K-Q-D-as} 
        \end{align}
       We want to show that we then also have:
       \[ \Kcal[A]:= \lB \bigotimes_{D \in \Dst[A]}^> \rB  \Kcal[D] = \Qf[A] \qquad \mu_{V\sm A}\text{-a.s.} \]
       For this fix a node $v \in D$ and note that 
        $Q(X_v|X_{\MBl^{G'}_<(v)})$ satisfies:
        \begin{align} P(X_{D_\le}|\doit(X_{J \cup V \sm D})) = Q(X_v|X_{\MBl^{G'}_<(v)}) \otimes 
            P(X_{D_<}|\doit(X_{J \cup V \sm D})). \label{eq:q-cond}
        \end{align}
        We then get the equalities:
            \begin{align*}
               & \tilde P(X_{D_\le}|\doit(X_{J \cup V \sm D})) \otimes \mu_{V\sm D}(X_{V\sm D}) \\
               &\stackrel{\text{Eq.\ }\ref{eq:K-Q-D-as}}{=}  P(X_{D_\le}|\doit(X_{J \cup V \sm D})) \otimes \mu_{V\sm D}(X_{V\sm D}) \\
               &\stackrel{\text{Eq.\ }\ref{eq:q-cond}}{=}  Q(X_v|X_{\MBl^{G'}_<(v)}) \otimes  P(X_{D_<}|\doit(X_{J \cup V \sm D})) \otimes \mu_{V\sm D}(X_{V\sm D}) \\ 
               &\stackrel{\text{Eq.\ }\ref{eq:K-Q-D-as}}{=}  Q(X_v|X_{\MBl^{G'}_<(v)}) \otimes  \tilde P(X_{D_<}|\doit(X_{J \cup V \sm D})) \otimes \mu_{V\sm D}(X_{V\sm D}). 
            \end{align*}
             So $Q(X_v|X_{\MBl^{G'}_<(v)})$ is a conditional of:
    \[  \tilde P(X_{D_\le}|\doit(X_{J \cup V \sm D})) \otimes \mu_{V\sm D}(X_{V\sm D}) \qquad \text{w.r.t.} \qquad
         \tilde P(X_{D_<}|\doit(X_{J \cup V \sm D}))\otimes \mu_{V\sm D}(X_{V\sm D}), \] 
         that  does not depend on $X_{A_\ge}$ (as $\MBl^{G'}_<(v) \ins A_<$).

      Now consider any other version of the conditional of:
         \[  \tilde P(X_{D_\le}|\doit(X_{J \cup V \sm D})) \otimes \mu_{V\sm D}(X_{V\sm D}) \qquad \text{w.r.t.} \qquad
          \tilde P(X_{D_<}|\doit(X_{J \cup V \sm D}))\otimes \mu_{V\sm D}(X_{V\sm D}), \] 
    that does not depend on variables attached to $A_\ge$ and which we will denote by:
        \[ K(X_v|X_{J \cup V \sm A_\ge}). \]
      Note that such a Markov kernel exists, as $Q(X_v|X_{\MBl^{G'}_<(v)})$ is such one.

      The same argumentation with $K(X_v|X_{J \cup V \sm A_\ge})$ in place of $Q(X_v|X_{\MBl^{G'}_<(v)})$, 
      using Eq.\ \ref{eq:K-Q-D-as}, shows that both, $Q(X_v|X_{\MBl^{G'}_<(v)})$ and $K(X_v|X_{J \cup V \sm A_\ge})$, 
    are then conditionals of
    \[  P(X_{D_\le}|\doit(X_{J \cup V \sm D})) \otimes \mu_{V\sm D}(X_{V\sm D}) \qquad \text{w.r.t.} \qquad
          P(X_{D_<}|\doit(X_{J \cup V \sm D}))\otimes \mu_{V\sm D}(X_{V\sm D}), \] 
    that do not depend on $X_{A_\ge}$.
    Now let:
       \[ \tilde N := \lC x_{J \cup V \sm A_\ge} \in \Xcal_{J \cup V \sm A_\ge}  
            \st K(X_v|X_{J \cup V \sm A_\ge}=x_{J \cup V \sm A_\ge})
        \neq Q(X_v|X_{\MBl^{G'}_<(v)} = x_{\MBl^G_<(v)}) \rC, \]
        and $N:= \tilde N \times \Xcal_{A_\ge} \ins \Xcal_{J \cup V}$. 
        Again, since both are versions of the same conditional we get:
        \[ \lp \Qf[D] \otimes \mu_{V \sm D} \rp (N^{x_J}|x_J) = 0,\]
        for every $x_J \in \Xcal_J$. Since $\Qf[D] \ll \mu_D$ we get for every $x_J \in \Xcal_J$:
        \[ \lp \mu_A \otimes \mu_{V \sm A} \rp (N^{x_J}) = \mu_V (N^{x_J}) =\lp \mu_D \otimes \mu_{V \sm D} \rp (N^{x_J}) = 0.\]
        Since also $\Qf[A] \ll \mu_A$ we get for every $x_J \in \Xcal_J$:
            \[ \lp \Qf[A] \otimes \mu_{V \sm A} \rp (N^{x_J}|x_J) = 0.\]
        Since $N= \tilde N \times \Xcal_{A_\ge}$ we get for every $x_J \in \Xcal_J$:
        \[ \lp P(X_{A_<}|\doit(X_{J \cup V \sm A})) \otimes \mu_{V\sm A}(X_{V \sm A}) \rp (\tilde N^{x_J}|x_J) = 
        \lp \Qf[A] \otimes \mu_{V \sm A} \rp (N^{x_J}|x_J) = 0.  \]
       This shows that the Markov kernels $Q(X_v|X_{\MBl^{G'}_<(v)})$ and $K(X_v|X_{J \cup V \sm A_\ge})$ are equal up to some
        $\lp P(X_{A_<}|\doit(X_{J \cup V \sm A})) \otimes \mu_{V\sm A}(X_{V \sm A}) \rp$-null set.
        Note that with this we get the factorization, using $A=\lC v_1,\dots,v_K \rC$, $v_1 < \cdots < v_K$:
        \begin{align*}
           & \Qf[A] \otimes \mu_{V \sm A} \\
           & = P(X_{A_{\le v_K}}|\doit(X_{J \cup V \sm A})) \otimes \mu_{V \sm A}(X_{V \sm A}) \\
           & = Q(X_{v_K}|X_{\MBl^{G'}_<(v_K)}) \otimes P(X_{A_{< v_K}}|\doit(X_{J \cup V \sm A})) \otimes \mu_{V \sm A}(X_{V \sm A}) \\
           & = K(X_{v_K}|X_{J \cup V \sm A_{\ge v_K}}) \otimes P(X_{A_{< v_K}}|\doit(X_{J \cup V \sm A})) \otimes \mu_{V \sm A}(X_{V \sm A}) \\
           & = K(X_{v_K}|X_{J \cup V \sm A_{\ge v_K}}) \otimes P(X_{A_{\le v_{K-1}}}|\doit(X_{J \cup V \sm A})) \otimes \mu_{V \sm A}(X_{V \sm A}) \\
           & = K(X_{v_K}|X_{J \cup V \sm A_{\ge v_K}}) \otimes Q(X_{v_K}|X_{\MBl^{G'}_<(v_K)}) \otimes P(X_{A_{< v_{K-1}}}|\doit(X_{J \cup V \sm A})) \otimes \mu_{V \sm A}(X_{V \sm A}) \\
           & = \cdots \\
           & = \lp \bigotimes_{v \in A}^> K(X_v|X_{J \cup V \sm A_\ge}) \rp \otimes \mu_{V \sm A}(X_{V \sm A})\\
           &=  \lp \lB \bigotimes_{D \in \Dst[A]}^> \rB \Kcal[D] \rp\otimes \mu_{V \sm A}(X_{V \sm A})\\
           &= \Kcal[A] \otimes \mu_{V \sm A}(X_{V \sm A}).
        \end{align*}
        Since such factorizations are essentially unique we get:
               \[ \Kcal[A] = \Qf[A] \qquad \mu_{V\sm A}\text{-a.s.} \]
        This shows the claim.

       c) Now let $D \ins V$ and $A \ins D$ with $A=\Anc^{[D]}(A)$. Consider that we are given a Markov kernel:
       \[ \Kcal[D]=\tilde P(X_{D}|\doit(X_{J \cup V \sm D})) :\, \Xcal_{J \cup V \sm D} \dshto \Xcal_{D}, \]
        such that:
       \[ \Kcal[D] = \Qf[D]\qquad \mu_{V\sm D}\text{-a.s.,}  \]
       which implies the equality:
            \[  \Kcal[D] \otimes \mu_{V\sm D} = \Qf[D]\otimes \mu_{V\sm D}. \]
            We want to show that the $A$-marginal of $\Kcal[D]$ equals $\Qf[A]$ up to $\mu_{V\sm A}$-null set.

        For this let $\Kcal[A]$ be the $A$-marginal of $\Kcal[D]$:
       \[ \Kcal[A]:=\tilde P(X_A|\doit(X_{J \cup V \sm D})) :\, \Xcal_{J \cup V \sm A} \to \Xcal_{J \cup V \sm D} \dshto \Xcal_A. \]
        Note that $\Qf[A]$ is the $A$-marginal of $\Qf[D]$.
        Marginalizing out $X_{D \sm A}$ on both sides in the above equation gives us:
         \[ \Kcal[A] \otimes \mu_{V\sm D} = \Qf[A]\otimes \mu_{V\sm D}. \]
        Multiplying both sides with $\mu_{D \sm A}$ gives:
        \[ \Kcal[A] \otimes \mu_{V\sm A} = \Kcal[A] \otimes \mu_{V\sm D} \otimes \mu_{D \sm A} = \Qf[A]\otimes \mu_{V\sm D} \otimes \mu_{D \sm A} = \Qf[A] \otimes \mu_{V\sm A}. \]
        Since such factorizations are essentially unique we get:
               \[ \Kcal[A] = \Qf[A] \qquad \mu_{V\sm A}\text{-a.s.} \]
        This shows the claim.

        This covers all cases of the ID-algorithm and thus shows the claim.
    \end{proof}
\end{Thm}


\begin{Thm}[Soundness up to continuous choices for strictly positive CBNs]
    \label{thm:ID-soundness-cont}
    Let $G=(J,V,E,L)$ be a CADMG with a fixed topological order $<$. % and $B \ins V$, $C \ins J \cup V$ two disjoint subsets.
    Consider the class of IL-CBNs $M$ with observational CADMG $G$ such that the following holds: 
    \begin{enumerate}
        \item The spaces $\Xcal_v$ are Polish spaces for $v \in J \cup V$,
        \item for every subset $D \ins V$ the interventional Markov kernel $\Qf[D]=P(X_D|\doit(X_{J \cup V \sm D}))$ 
            is \emph{strictly positive} (on non-empty open subsets of $\Xcal_{D}$), and:
        \item for every $v \in D$ the Markov kernel $Q(X_v|X_{\MBl^{G_{\doit(D^\cmpl)}}_<(v)})$ 
            can be chosen to be \emph{continuous}, viewed as a map: $\Xcal_{\MBl^{G_{\doit(D^\cmpl)}}_<(v)} \to \Prob(\Xcal_v)$.
    \end{enumerate}
    If the ID-algorithm does not produce FAIL for input $B, C \ins G$,
    then $P(X_B|\doit(X_{J \cup C}))$ is identifiable
    and \trackable\ ``up to continuous choices of conditional Markov kernels'' 
    from $P(X_V|\doit(X_J))$ and $G$ for such CBNs $M$, 
    i.e.\ if every occuring conditional Markov kernel is chosen to be continuous (which will always be possible by the assumptions made).
    \begin{proof}
        For a district $D \in \Dst[V]$ and $v \in D$, 
        by assumption, there exists a \emph{continuous} version of $Q(X_v|X_{\MBl^G_<(v)})$, which is also a version of:
     \[ P(X_v|X_{\Pred^{[V]}_<(v)},\doit(X_J)) \qquad \text{and} \qquad P(X_v|X_{\Pred^{[D]}_<(v)},\doit(X_{J \cup V\sm D})).\]
      We abbreviate $V_<:=\Pred^{[V]}_<(v)$ and $D_<:=\Pred^{[D]}_<(v)$ in the following.

     Now consider any \emph{continuous} version of the conditional Markov kernel $P(X_v|X_{V_<},\doit(X_J))$.
     Note that such a version always exists because $Q(X_v|X_{\MBl^G_<(v)})$ is already an existing continuous version.
     Since $P(X_{V_<}|\doit(X_J))$ is strictly positive, as the marginal of $\Qf[V]$,
      Lemma \ref{lem:cont-cond-unique} implies then the ``sure'' equality:
     \[ P(X_v|X_{V_<},\doit(X_J)=Q(X_v|X_{\MBl^G_<(v)}).  \]
     So any continuous version of $P(X_v|X_{V_<},\doit(X_J)$ necessarily agrees with $Q(X_v|X_{\MBl^G_<(v)})$ on all points.

     Similarly, using the same arguments, we get that $Q(X_v|X_{\MBl^G_<(v)})$ is ``surely'' equal to every \emph{continuous} version 
     of the conditional
     $P(X_v|X_{D_<},\doit(X_{J \cup V\sm D}))$, as also the Markov kernel $P(X_{D_<}|\doit(X_{J \cup V\sm D}))$ is strictly positive,
     as a marginal of $\Qf[D]$. So we get:
     \[ Q(X_v|X_{\MBl^G_<(v)}) = P(X_v|X_{D_<},\doit(X_{J \cup V\sm D})).  \]

     This means that if we pick a/the \emph{continuous} version of the conditional $P(X_v|X_{V_<},\doit(X_J)$ 
     then it is ``surely'' equal to
     the/every \emph{continuous} version of the conditional $P(X_v|X_{D_<},\doit(X_{J \cup V\sm D}))$:
     \[ P(X_v|X_{V_<},\doit(X_J)=Q(X_v|X_{\MBl^G_<(v)}) = P(X_v|X_{D_<},\doit(X_{J \cup V\sm D})).  \]

     These arguments, repeated for subgraphs, then show that in the ID-algorithm for every occuring conditional and product
     (e.g.\ for districts and the final product) we end up with 
     distinct and correct choices for all Markov kernels. This then also shows the identifiability of such CBNs $M$
     (in the not-FAIL case).
    \end{proof}
\end{Thm}


\begin{Lem}
    \label{lem:cbn-str-pos-density}
    Consider an IL-BCN:
        \[M= \lp G^+=\lp J,(V,U),E^+ \rp, \lp P_v ( X_v|\doit( X_{\Pa^{G^+}(v)})) \rp_{v \in V \cup U} \rp,\]
        with observational CADMG $G=(J,V,E,L)$ and fixed topological order $<$.
    Assume that for every $v \in V$ we have a measure $\mu_v$ on $\Xcal_v$ such that $P_v ( X_v|\doit( X_{\Pa^{G^+}(v)}))$
    has a (strictly positive) density w.r.t.\ $\mu_v$:
    \[ p(x_v|\doit(x_{\Pa^{G^+}(v)})) > 0. \]
    Furthermore, we put for $x_V \in \Xcal_V$, $x_U \in \Xcal_U$, $x_J \in \Xcal_J$:
    \[ p(x_V|x_U,\doit(x_J)) := \prod_{v \in V} p(x_v|\doit(x_{\Pa^{G^+}(v)})),  \]
    and then integrate in reverse order of $<$:
    \[ p(x_V|\doit(x_J)) := \idotsint_{\Xcal_U} p(x_V|x_U,\doit(x_J)) \, 
    \bigotimes_{ u \in U}^> P_u (X_u \in dx_u|\doit( X_{\Pa^{G^+}(v)} = x_{\Pa^{G^+}(v)} )).   \]
    Then the former is a (strictly positive) density of $P(X_V,X_U|\doit(X_J))$ w.r.t.\ 
    \[ \bigotimes_{ v \in V}^> \mu_v \otimes^> \bigotimes_{ u \in U}^> P_u (X_u|\doit( X_{\Pa^{G^+}(v)} )), \]
    and the latter a (strictly positive) density of $P(X_V|\doit(X_J))$ w.r.t.\ 
    \[ \mu_V := \bigotimes_{ v \in V} \mu_v. \]
    Similarly, for every $D \ins V$ the interventional Markov kernel $P(X_D|\doit(X_{J \cup V \sm D}))$ has
    a (strictly positive) density w.r.t.\ $\mu_D:=\bigotimes_{ v \in D} \mu_v$.
    \begin{proof}
    The claim can be shown by integrating the above densities over product sets $A = \prod_v A_v$.
    Inductively we can use Fubini's theorem and:
    \[ \int_{A_v} p(x_v|\doit(x_{\Pa^{G^+}(v)}))\, \mu_v(dx_v) = P_v(X_v \in A_v|\doit( X_{\Pa^{G^+}(v)}=x_{\Pa^{G^+}(v)})).\]
     Regarding strict positivity, note that if $f(x) > 0$ for all $x$, then $\int f\, d\mu >0$ for non-trivial $\mu$. 
        So strict positivity is preserved through integration.
    \end{proof}
\end{Lem}

